% =====================================================================
% CPE 503 — GROUP 09 — PROBLEM 9
% Smart Grid FDI Anomaly Detection Engine via CUSUM and Mahalanobis Distance
% Comprehensive Single-File Technical Report (Overleaf-Ready)
% Department of Computer Engineering, Redeemer's University, Ede
% =====================================================================

\documentclass[11pt,a4paper]{report}

% ---------------- PACKAGES ----------------
\usepackage[utf8]{inputenc}
\usepackage[T1]{fontenc}
\usepackage{lmodern}
\usepackage{geometry}
\usepackage{setspace}
\usepackage{amsmath,amssymb,amsfonts}
\usepackage{graphicx}
\usepackage{float}
\usepackage{placeins}
\usepackage{booktabs}
\usepackage{array}
\usepackage{multirow}
\usepackage{caption}
\usepackage{cite}
\usepackage{url}
\usepackage{listings}
\usepackage{xcolor}
\usepackage{fancyhdr}

% ---------------- PAGE FORMATTING ----------------
\geometry{left=0.85in,right=0.85in,top=0.85in,bottom=0.85in}
\setstretch{1.08}
\setlength{\headheight}{15pt}
\setlength{\parindent}{0.35in}
\setlength{\parskip}{0pt}
% ---------------- APA 7TH EDITION UNNUMBERED REFERENCES ----------------
\renewcommand{\bibname}{References}
\makeatletter
\def\@cite#1#2{({#1\if@tempswa , #2\fi})}
\def\@biblabel#1{}
\renewenvironment{thebibliography}[1]{%
  \chapter*{\bibname}%
  \addcontentsline{toc}{chapter}{\bibname}%
  \list{}{%
    \setlength{\labelwidth}{0pt}%
    \setlength{\labelsep}{0pt}%
    \setlength{\leftmargin}{0.35in}%
    \setlength{\itemindent}{-0.35in}%
    \setlength{\itemsep}{6pt}%
    \setlength{\parsep}{0pt}%
    \usecounter{enumiv}%
    \let\p@enumiv\@empty
    \renewcommand\theenumiv{\@arabic\c@enumiv}%
  }%
  \sloppy\clubpenalty4000\widowpenalty4000%
  \sfcode`\.=1000\relax
}{%
  \endlist
}
\makeatother




\pagestyle{fancy}
\fancyhf{}
\fancyhead[L]{Group 9}
\fancyhead[R]{Smart Grid FDI Anomaly Engine}
\fancyfoot[C]{\thepage}
\renewcommand{\headrulewidth}{0.4pt}

\graphicspath{{results/}}

\lstset{
  basicstyle=\ttfamily\footnotesize,
  breaklines=true,
  frame=single,
  keywordstyle=\color{blue},
  commentstyle=\color{gray},
  stringstyle=\color{teal},
  showstringspaces=false,
  numbers=left,
  numberstyle=\tiny,
  xleftmargin=1.5em,
  framexleftmargin=1em
}

% ---------------- SAFE FIGURE COMMAND (SUPPORTS PNG & JPG) ----------------
\newcommand{\reportfigure}[4][0.58\textwidth]{%
\begin{figure}[htbp]
\centering
\IfFileExists{results/#2}{%
  \includegraphics[width=#1]{results/#2}%
}{%
  \IfFileExists{results/#2.png}{%
    \includegraphics[width=#1]{results/#2.png}%
  }{%
    \IfFileExists{results/#2.jpg}{%
      \includegraphics[width=#1]{results/#2.jpg}%
    }{%
      \fbox{\parbox[c][5cm][c]{0.82\textwidth}{\centering
      \textbf{Figure file not supplied}\\[0.4em]
      Expected file: \texttt{results/#2} (or \texttt{.png} / \texttt{.jpg})\\[0.4em]
      Upload the corresponding generated plot to the Overleaf
      project before final submission.}}%
    }%
  }%
}
\caption{#3}
\label{#4}
\end{figure}
}

% ---------------- DOCUMENT INFORMATION ----------------
\title{\textbf{SMART GRID FALSE DATA INJECTION (FDI) ANOMALY DETECTION ENGINE VIA CUSUM AND MAHALANOBIS DISTANCE}}
\author{Group 9\\Department of Computer Engineering\\Redeemer's University, Ede}
\date{November 2026}

\begin{document}


% =====================================================================
% TITLE PAGE
% =====================================================================
\begin{titlepage}
\centering
\vspace*{1cm}
{\Large\textbf{REDEEMER'S UNIVERSITY, EDE}}\\[0.4cm]
{\large Department of Computer Engineering}\\
{\large Faculty of Engineering}

\vspace{1.5cm}
{\Large\textbf{GROUP 9 PROJECT REPORT}}\\[1cm]

{\LARGE\textbf{
SMART GRID FALSE DATA INJECTION (FDI)\\
ANOMALY DETECTION ENGINE VIA CUSUM AND\\
MAHALANOBIS DISTANCE
}}

\vspace{1.5cm}
\textbf{A Technical Report Submitted in Partial\\
Fulfilment of the Requirements for CPE 503\\
(Embedded Systems)}

\vspace{1.5cm}
\textbf{GROUP MEMBERS}\\[0.5cm]

\textsc{Adediwura Esther Oluwarasimi}\\
-- RUN/CPE/22/12005\\[0.3cm]

\textsc{Ololade John Ayotunde}\\
-- RUN/CPE/22/12072\\[0.3cm]

\textsc{Akubo Bright}\\
-- RUN/CPE/22/12025\\[0.3cm]

\textsc{Ojo Oluwadamilare Samuel}\\
-- RUN/CPE/22/12062

\vfill
\textbf{November 2026}
\end{titlepage}

% =====================================================================
% PRELIMINARY PAGES
% =====================================================================
\pagenumbering{roman}

\chapter*{Declaration and Acknowledgements}
\addcontentsline{toc}{chapter}{Declaration and Acknowledgements}

\subsection*{Declaration of Originality}
We declare that this report presents the original work undertaken by Group 9
for the CPE 503 (Embedded Systems) project in the Department of Computer
Engineering, Redeemer's University, Ede. All sources consulted, including
academic literature, benchmarks, and software libraries, have been duly
acknowledged and referenced. The numerical results, statistical metrics,
and hardware resource estimates presented herein reflect authentic simulation
and synthesis executions conducted on the project codebase. In accordance with
the strict academic standards of the course, no unperformed hardware test is
represented as a completed physical experiment, and all design boundaries,
unconstrained timing paths, and interface limitations are reported with complete
academic integrity.

\vspace{0.8cm}
\subsection*{Acknowledgements}
The members of Group 9 express their sincere gratitude to the course
lecturer, the Department of Computer Engineering, and the Faculty of
Engineering at Redeemer's University, Ede, for their academic leadership,
rigorous standards, and the opportunity to undertake this research project.
We also acknowledge the open-source community, IEEE, and the developers of
scientific and power-system computing tools (Python, NumPy, SciPy, Matplotlib,
and MATPOWER) as well as the AMD/Xilinx Vivado Design Suite for providing the
computational infrastructure that supported our modelling and RTL synthesis.

\chapter*{Abstract}
\addcontentsline{toc}{chapter}{Abstract}
False Data Injection (FDI) attacks constitute an existential threat to
modern smart electrical grids. By manipulating sensor measurements
transmitted across Wide Area Monitoring Systems (WAMS), an adversary can
deceive Power System State Estimation (PSSE) algorithms into producing
erroneous operating states. Coordinated stealth attacks, constructed along
the column space of the power-flow measurement Jacobian ($\boldsymbol{a} = \boldsymbol{H}\boldsymbol{c}$),
leave the measurement residual strictly invariant, rendering classical
chi-square ($\chi^2$) Bad Data Detection (BDD) mathematically blind.

To overcome this fundamental vulnerability, this project presents the
design, mathematical formulation, and hardware architecture of a dual-layer
statistical anomaly detection engine implemented in SystemVerilog for the
AMD/Xilinx Artix-7 XC7A200T FPGA. Layer 1 executes an instantaneous,
covariance-aware Mahalanobis distance test that exploits physical cross-channel
correlations among synchrophasor measurements ($V_1, V_2, P_{12}, Q_{12}$),
while Layer 2 implements a recursive, two-sided Cumulative Sum (CUSUM) sequential
test to capture subtle, persistent parameter drifts that evade single-frame
thresholds.

A golden reference model in Python evaluates the detector across 1,200 continuous
time-series observations derived from an IEEE 14-bus transmission system,
subjected to four attack modalities: coordinated stealth FDI, gradual ramp drift,
false scaling, and historical replay. Across the evaluated benchmark, the proposed
hybrid engine achieves 84.67\% accuracy, 96.36\% precision, 71.57\% recall, an
F1-score of 0.8214, and an ROC-AUC of 0.941, while cutting detection latency for
slow ramp attacks by 4.2$\times$ compared to conventional BDD.

The synthesizable architecture employs a signed Q1.15 fixed-point datapath
featuring 34-bit accumulator guard bits, convergent half-up rounding, and
saturation logic to guarantee numerical stability. Vivado 2024.1 synthesis for
the Artix-7 target utilizes 1,193 Slice LUTs (0.89\%), 647 Slice Registers
(0.24\%), 24 DSP48E1 slices (3.24\%), and 2 RAMB18E1 primitives (0.27\%). At a
100~MHz target clock, the fully pipelined engine processes a 4-channel measurement
vector in 36 clock cycles (360~ns), providing a maximum throughput of 2.78
MSamples/s and confirming the architectural viability of line-rate anomaly
detection for modern transmission substations.

\vspace{0.8cm}
\noindent\textbf{Keywords:} Smart Grid Cybersecurity, False Data Injection (FDI),
Mahalanobis Distance, Two-Sided CUSUM, FPGA Architecture, SystemVerilog, Artix-7.

\clearpage
\tableofcontents
\listoffigures
\listoftables
\clearpage
\pagenumbering{arabic}


% =====================================================================
% CHAPTER ONE
% =====================================================================
\chapter{Introduction}

\section{Background of the Study}
Modern electrical power transmission and distribution networks are rapidly
evolving from isolated, mechanically switched systems into highly automated,
interconnected cyber-physical infrastructures commonly referred to as smart
grids. This transformation is driven by the imperative to enhance grid
reliability, dynamically integrate high-penetration renewable energy sources
(such as photovoltaic arrays and wind generation governed by standards such as IEEE~1547 \cite{ieee1547}), and optimize economic power
dispatch. To maintain situational awareness across geographically dispersed
substations, smart grids rely on sophisticated Wide Area Monitoring Systems
(WAMS), Supervisory Control and Data Acquisition (SCADA) networks, and
high-speed synchrophasor measurement technology.

Phasor Measurement Units (PMUs) form the technological backbone of modern
transmission monitoring. Operating in accordance with standards such as
IEEE~C37.118 and its dual-logo successor IEC/IEEE~60255-118-1 \cite{ieeeC37118,iec60255},
PMUs sample analog voltage and current waveforms synchronously using Global
Positioning System (GPS) radio clocks, reporting time-stamped positive-sequence
phasor estimates at reporting rates typically ranging between 50 and 60 frames
per second ($F_s = 50\text{--}60\text{ Hz}$). These high-cadence data streams
are transmitted over utility communication networks to control-room Energy
Management Systems (EMS), where Power System State Estimation (PSSE) algorithms
filter measurement noise, reconcile topological inconsistencies, and estimate
the complex voltage vector across all network buses.

However, the deep convergence of operational technology (OT) with standard
information and communication technology (ICT) protocols introduces an extensive
cyber-attack surface. Transmission telemetry traverses microwave links, fiber-optic
substation local area networks, and wide-area routed backbones that are vulnerable
to packet interception, physical tapping, and man-in-the-middle (MITM) manipulation.
Among the various cyber threats facing modern power grids, False Data Injection
(FDI) attacks represent one of the most stealthy and damaging attack classes.
By systematically corrupting sensor telemetry, an adversary can deliberately
mislead state estimation engines, inducing operators or automated protection
mechanisms to execute incorrect control actions, such as unwarranted generator
tripping, line disconnections, or malicious power redispatch, ultimately risking
cascading system-wide blackouts.

\section{Statement of the Problem}
Power System State Estimation typically relies on Weighted Least-Squares (WLS)
formulations to extract the most probable system state $\boldsymbol{x} \in \mathbb{R}^n$
from a redundant set of noisy measurements $\boldsymbol{z} \in \mathbb{R}^m$, where
$m > n$. To protect against sensor faults and telemetry corruption, conventional
EMS deployments implement Bad Data Detection (BDD) algorithms based on residual
analysis. Specifically, the measurement residual vector is defined as
\begin{equation}
\boldsymbol{r} = \boldsymbol{z} - \boldsymbol{H}\hat{\boldsymbol{x}},
\label{eq:ch1_residual}
\end{equation}
where $\boldsymbol{H} \in \mathbb{R}^{m \times n}$ is the power-system measurement
Jacobian and $\hat{\boldsymbol{x}}$ is the calculated state estimate. Under nominal
conditions with zero-mean Gaussian noise, the squared weighted residual norm
$J(\hat{\boldsymbol{x}}) = \boldsymbol{r}^T \boldsymbol{R}^{-1} \boldsymbol{r}$
follows a central chi-square distribution $\chi^2_{(m-n)}$. If $J(\hat{\boldsymbol{x}})$
exceeds a predefined significance threshold $\gamma_{\chi^2}$, an alarm is asserted,
and bad data identification procedures are initiated.

However, seminal research by Liu, Ning, and Reiter \cite{liu2009,liu2011} demonstrated
a profound mathematical vulnerability in conventional BDD. If an adversary acquires
sufficient knowledge of the power network topology and line parameters to construct
the measurement Jacobian $\boldsymbol{H}$, the attacker can craft an injection
vector:
\begin{equation}
\boldsymbol{a} = \boldsymbol{H}\boldsymbol{c},
\label{eq:ch1_stealth}
\end{equation}
where $\boldsymbol{c} \in \mathbb{R}^n$ is an arbitrary false state bias vector. When
this perturbation is added to the clean telemetry stream ($\boldsymbol{z}_a = \boldsymbol{z} + \boldsymbol{a}$),
the resulting attacked state estimate shifts precisely by $\hat{\boldsymbol{x}}_a = \hat{\boldsymbol{x}} + \boldsymbol{c}$.
Consequently, the post-attack measurement residual is
\begin{equation}
\boldsymbol{r}_a = \boldsymbol{z}_a - \boldsymbol{H}\hat{\boldsymbol{x}}_a = (\boldsymbol{z} + \boldsymbol{H}\boldsymbol{c}) - \boldsymbol{H}(\hat{\boldsymbol{x}} + \boldsymbol{c}) = \boldsymbol{z} - \boldsymbol{H}\hat{\boldsymbol{x}} = \boldsymbol{r}.
\end{equation}
Because the post-attack residual $\boldsymbol{r}_a$ is \textbf{strictly identical}
to the pre-attack residual $\boldsymbol{r}$, the chi-square detection statistic
remains completely unchanged ($J(\hat{\boldsymbol{x}}_a) = J(\hat{\boldsymbol{x}})$).
Conventional residual-based bad data detection is therefore mathematically blind to
stealth FDI attacks.

Furthermore, existing software-based machine learning and statistical detection
schemes implemented on general-purpose CPUs or GPUs introduce non-deterministic
operating system scheduling jitter and inference latencies exceeding tens of
milliseconds, rendering them unsuitable for real-time, line-rate deployment inside
substation Intelligent Electronic Devices (IEDs) or streaming Phasor Data Concentrators
(PDCs). An urgent engineering need therefore exists for dedicated, hardware-accelerated
statistical anomaly detection engines capable of examining both multivariate spatial
correlations and sequential temporal persistence within sub-microsecond latency bounds.

\section{Aim of the Project}
The primary aim of this project is to design, implement, and experimentally verify
a high-throughput, hardware-oriented statistical anomaly detection engine combining
covariance-aware Mahalanobis distance and recursive two-sided Cumulative Sum (CUSUM)
filtering on an AMD/Xilinx Artix-7 FPGA, providing robust detection of stealth and
dynamic False Data Injection attacks on smart-grid measurement streams.

\section{Specific Objectives}
To achieve the stated aim, the project is executed through the following measurable
and testable specific objectives:
\begin{enumerate}
    \item \textbf{Primary Objective:} To design, code, and synthesize a fully
    pipelined, synthesizable SystemVerilog (IEEE~1800-2017) hardware core on an
    AMD/Xilinx Artix-7 XC7A200T-FBG676-2 FPGA that processes incoming multi-channel
    PMU telemetry within a guaranteed deterministic latency of under $1.0~\mu\text{s}$
    per 4-channel vector at a 100~MHz clock target, utilizing less than 5\% of
    chip resources.
    \item \textbf{Mathematical Modelling Objective:} To develop a bit-accurate
    Python golden reference model that formulates the joint spatial Mahalanobis
    distance and two-sided CUSUM sequential test, deriving closed-form expressions
    for detection thresholds ($\tau^2, h$) and reference slack ($k$) based on
    statistical significance levels ($\alpha$).
    \item \textbf{Power Grid Telemetry Synthesis:} To generate a comprehensive
    1,200-frame time-series benchmark dataset derived from an IEEE 14-bus transmission
    system, simulating four distinct cyber-attack regimes: coordinated stealth FDI
    ($\boldsymbol{a} = \boldsymbol{H}\boldsymbol{c}$), gradual ramp drift,
    false scaling, and historical replay.
    \item \textbf{Verification and Comparative Evaluation:} To conduct extensive
    statistical and hardware verification demonstrating zero bit-level discrepancy
    between the SystemVerilog RTL simulation and the Python golden model, and to
    benchmark detection performance (Accuracy, Precision, Recall, F1-score, ROC-AUC,
    and detection delay) against classical $\chi^2$ bad data detection.
\end{enumerate}

\section{Research Questions}
This investigation is guided by two fundamental engineering research questions:
\begin{itemize}
    \item \textbf{Research Question 1 (RQ1):} To what extent does the architectural
    coupling of an instantaneous spatial Mahalanobis distance metric with a temporal
    two-sided CUSUM filter mitigate the residual blindness of classical state
    estimation, and what are the quantitative performance trade-offs across abrupt
    versus gradual cyber-attack modalities?
    \item \textbf{Research Question 2 (RQ2):} Can a signed Q1.15 fixed-point
    arithmetic datapath incorporating 34-bit accumulator guard bits, convergent
    half-up rounding, and overflow saturation logic preserve statistical detection
    accuracy within 1\% of double-precision floating-point software models while
    meeting strict line-rate throughput constraints on an Artix-7 FPGA?
\end{itemize}

\section{Stated Hypotheses and Success Criteria}
Prior to running experimental simulations, the following technical hypotheses and
success criteria were established:
\begin{itemize}
    \item \textbf{Hypothesis 1 (Detection Superiority):} Because coordinated stealth
    injections ($\boldsymbol{a} = \boldsymbol{H}\boldsymbol{c}$) induce cross-channel
    power-flow and voltage violations despite preserving the measurement residual,
    the proposed dual-layer detector will achieve a detection recall exceeding 70\%
    and an F1-score $> 0.80$, compared to $< 25\%$ recall for classical $\chi^2$ BDD.
    \item \textbf{Hypothesis 2 (False Alarm Suppression):} Setting the CUSUM reference
    allowance to half the design shift magnitude ($k = \delta/2$) will suppress
    background Gaussian noise fluctuations, achieving a false positive rate (FPR)
    below 3.0\% during clean, unattacked operational windows.
    \item \textbf{Criterion 3 (Hardware Determinism):} The SystemVerilog RTL core
    must exhibit bit-exact mathematical concordance (100\% known-answer test pass
    rate) against the Python fixed-point golden model across all 1,200 simulation frames.
\end{itemize}

\section{Key Contributions of the Project}
The principal contributions delivered by Group 9 in this work are summarized as follows:
\begin{enumerate}
    \item \textbf{Dual-Layer Spatial-Temporal Architecture:} A unified anomaly detection
    framework that synergistically combines spatial multivariate correlation (Layer 1:
    Mahalanobis distance) for immediate identification of abrupt structural shifts with
    temporal sequential accumulation (Layer 2: two-sided CUSUM) for detecting slow,
    sub-threshold ramp drifts.
    \item \textbf{Robust Fixed-Point Arithmetic Datapath:} A fully synthesizable,
    parameterized Q1.15 RTL architecture engineered with 34-bit accumulator guard bits
    to eliminate dot-product overflow, an unbiased $+2^{14}$ half-up rounding unit, and
    two's-complement saturation logic.
    \item \textbf{Comprehensive Multi-Attack Evaluation Suite:} An extensive experimental
    benchmark across four challenging cyber-attack classes on IEEE 14-bus telemetry,
    providing empirical evidence of detection delays, Average Run Length ($ARL_0/ARL_1$)
    trade-offs, and Receiver Operating Characteristic (ROC) curves.
    \item \textbf{Open, Reproducible Hardware/Software Repository:} A self-contained,
    verifiable artifact package containing SystemVerilog source code, self-checking
    testbenches, Python golden models, Vivado synthesis Tcl scripts, raw CSV telemetry,
    and plotting pipelines.
\end{enumerate}

\section{Scope and Limitations of the Project}
The scope of this project encompasses the algorithm design, mathematical derivation,
fixed-point error analysis, SystemVerilog RTL coding, Vivado out-of-context logic
synthesis, and Python-based statistical evaluation targeting an AMD/Xilinx Artix-7
XC7A200T-FBG676-2 FPGA (as featured on the AC701 evaluation board). Telemetry streams
are generated from a 4-channel measurement configuration ($V_1, V_2, P_{12}, Q_{12}$)
derived from the IEEE 14-bus transmission benchmark.

In strict adherence to academic honesty, the following project boundaries are declared:
\begin{itemize}
    \item The reported hardware metrics reflect out-of-context and top-level logic
    synthesis using Vivado 2024.1; physical bitstream generation and in-system
    validation on physical AC701 hardware boards were not executed due to laboratory
    bench constraints.
    \item The raw top-level entity exposes 572 parallel I/O pins, exceeding the 400
    bonded IOBs available on the XC7A200T package. Consequently, board-level deployment
    mandates an AXI4-Stream or UART serial wrapper.
    \item The statistical model assumes stationary background Gaussian noise and a
    precomputed inverse covariance matrix; adaptive tracking of topological switching
    transients remains an area for future work.
\end{itemize}

\section{Organisation of the Report}
The remainder of this report is organized as follows:
\begin{itemize}
    \item \textbf{Chapter Two (Literature Review):} Provides a structured survey of
    smart grid architecture, cybersecurity vulnerabilities, FDI attack mechanics,
    conventional BDD limitations, Mahalanobis and CUSUM theory, FPGA implementations,
    a comparative analysis of published works, and the explicit research gap.
    \item \textbf{Chapter Three (Mathematical Formulation and Fixed-Point Design):}
    Presents the formal power-system state estimation equations, mathematical proofs
    of stealth FDI evasion, closed-form Mahalanobis and CUSUM formulations, fixed-point
    Q1.15 scaling bounds, and hardware complexity models.
    \item \textbf{Chapter Four (Methodology and Hardware Implementation):} Details the
    10-step design methodology, Python telemetry generation, SystemVerilog datapath
    design, three-process Moore controller FSM, and verification hierarchy.
    \item \textbf{Chapter Five (Results and Discussion):} Examines the empirical
    detection performance, FPGA resource utilization, timing and power estimates,
    multi-scenario ablation results, and 9 visual analytics plots.
    \item \textbf{Chapter Six (Conclusion and Recommendations):} Concludes the report
    by evaluating achievements against stated objectives and outlining practical next
    steps for physical system deployment.
\end{itemize}


% =====================================================================
% CHAPTER TWO
% =====================================================================
\chapter{Literature Review}

\section{Introduction}
Modern smart electrical grids integrate advanced communication, digital sensing, and automated
control platforms with traditional high-voltage physical assets. While Phasor Measurement Units
(PMUs) and Wide Area Monitoring Systems (WAMS) enable real-time situational awareness, their
reliance on networked communication introduces severe cybersecurity vulnerabilities. Among these,
False Data Injection (FDI) attacks deliberately manipulate sensor telemetry to deceive Power
System State Estimation (PSSE) while evading conventional bad-data detection.

To address this challenge, researchers have explored multivariate statistical analysis and
sequential change-point monitoring. Mahalanobis distance evaluates instantaneous cross-channel
correlations, while the Cumulative Sum (CUSUM) algorithm integrates subtle persistent shifts over
time. Translating these algorithms into digital hardware such as Field-Programmable Gate Arrays
(FPGAs) provides deterministic microsecond execution, but requires careful fixed-point scaling
to manage finite-precision effects. This chapter reviews the architectural, statistical, and
hardware foundations of smart-grid FDI anomaly detection, culminating in a critical analysis of
existing approaches and an explicit research gap.

\section{Smart Grid Architecture and Operation}

\subsection{Overview and Conceptual Domains}
A smart grid functions as an interconnected cyber-physical system where physical electricity flows
are continuously supervised by digital information networks. The National Institute of Standards
and Technology (NIST) conceptual model \cite{nist2014framework} identifies seven functional domains,
as summarized in Table~\ref{tab:smartgrid_domains}.

\begin{table}[htbp]
\centering
\caption{Major domains of the NIST smart grid conceptual model.}
\label{tab:smartgrid_domains}
\renewcommand{\arraystretch}{1.2}
\small
\begin{tabular}{|p{3.0cm}|p{5.0cm}|p{6.0cm}|}
\hline
\textbf{Domain} & \textbf{Primary Function} & \textbf{Typical Components or Activities} \\
\hline
Bulk Generation & Large-scale electricity generation & Power plants, renewable farms, generation control \\
\hline
Transmission & Long-distance bulk power transfer & High-voltage lines, substations, PMUs, relays \\
\hline
Distribution & Local electricity delivery to end users & Feeders, reclosers, transformers, microgrids \\
\hline
Operations & Network supervision and control & SCADA, Energy Management Systems (EMS), WAMS \\
\hline
Markets & Economic energy exchange & Trading platforms, aggregators, clearing houses \\
\hline
Service Providers & Third-party utility support & Billing, data analytics, customer services \\
\hline
Customers & Energy consumption and local generation & Smart meters, EV chargers, rooftop solar arrays \\
\hline
\end{tabular}
\end{table}

\subsection{Power Flows and Data Integrity}
Transmission networks transmit electric energy across geographically distributed buses. Monitoring
relies on voltage, current, and active/reactive power measurements to estimate network operating
states. Because supervisory decisions directly impact physical grid stability, measurement
integrity is paramount: corrupted telemetry can induce erroneous generation redispatch, line
tripping, or cascading outages.

\section{Cybersecurity in Smart Grid Systems}
In power system operational technology (OT), the classic CIA triad prioritizes \textbf{Availability}
and \textbf{Integrity} over \textbf{Confidentiality} \cite{nist80082}. Telemetry streams must arrive
deterministically at 50~Hz (every 20~ms) to drive protection routines. Attack vectors include
substation LAN intrusion, unencrypted protocol tampering (DNP3, Modbus, IEC 60870-5-104), and
GPS clock spoofing. While standards such as NISTIR 7628 \cite{nist7628,nist7628r1} and IEC 62443
prescribe perimeter defenses, secondary statistical anomaly detection engines inside substations
remain essential to counter internal or stealthy network breaches.

\section{False Data Injection Attacks}
Let nominal telemetry be $\boldsymbol{z} = \boldsymbol{H}\boldsymbol{x} + \boldsymbol{e}$. In an FDI
attack, the adversary injects an additive perturbation: $\boldsymbol{z}_a = \boldsymbol{z} + \boldsymbol{a}$.
As proven by Liu, Ning, and Reiter \cite{liu2009}, if the adversary constructs the attack along the
measurement Jacobian:
\begin{equation}
\boldsymbol{a} = \boldsymbol{H}\boldsymbol{c},
\end{equation}
where $\boldsymbol{c}$ is an arbitrary state bias, the resulting measurement residual $\boldsymbol{r}_a$
remains \textbf{strictly invariant}, rendering the attack unobservable to state estimation.
Depending on execution temporal profiles, FDI attacks manifest as abrupt step injections, gradual
ramps, multi-channel scaling, or historical replay. Crucially, detectability is not automatic: it
depends on monitored physical quantities, topological coupling, and injection magnitude relative to noise.

\section{Conventional Bad-Data Detection Techniques}
Conventional state estimation computes the Weighted Least-Squares (WLS) state estimate
$\hat{\boldsymbol{x}} = (\boldsymbol{H}^T\boldsymbol{R}^{-1}\boldsymbol{H})^{-1}\boldsymbol{H}^T\boldsymbol{R}^{-1}\boldsymbol{z}$
and evaluates the chi-square residual test $J(\hat{\boldsymbol{x}}) = \boldsymbol{r}^T\boldsymbol{R}^{-1}\boldsymbol{r} \lessgtr \gamma_{\chi^2}$
\cite{abur2004}. Under stealth FDI ($\boldsymbol{a} = \boldsymbol{H}\boldsymbol{c}$):
\begin{equation}
\boldsymbol{r}_a = \boldsymbol{z}_a - \boldsymbol{H}\hat{\boldsymbol{x}}_a = (\boldsymbol{z} + \boldsymbol{H}\boldsymbol{c}) - \boldsymbol{H}(\hat{\boldsymbol{x}} + \boldsymbol{c}) = \boldsymbol{r}.
\end{equation}
Because the post-attack residual $\boldsymbol{r}_a$ is mathematically identical to pre-attack
residual $\boldsymbol{r}$, classical $\chi^2$ and Largest Normalized Residual (LNR) tests fail
completely, establishing the necessity for alternative statistical features.

\section{Statistical Anomaly Detection}
Power grid variables ($V_1, V_2, P_{12}, Q_{12}$) are physically coupled by Kirchhoff's laws and line
impedance. Univariate limit checkers fail against coordinated attacks whose elements remain within
individual bounds. Multivariate monitoring evaluates joint probability density functions, exploiting
cross-channel correlations to flag anomalies. Effective defense requires addressing both spatial
cross-correlation (instantaneous coherence) and temporal persistence (evidence accumulation).

\section{Mahalanobis Distance}
The Mahalanobis distance \cite{mahalanobis1936} measures statistical divergence from a calibrated
baseline $(\boldsymbol{\mu}, \boldsymbol{\Sigma})$:
\begin{equation}
D^2 = (\boldsymbol{q} - \boldsymbol{\mu})^T \boldsymbol{\Sigma}^{-1} (\boldsymbol{q} - \boldsymbol{\mu}) \sim \chi^2_d.
\end{equation}
By rotating deviations into principal axes and weighting by inverse variance, $D^2$ instantly
detects coordinate-conforming injections ($\boldsymbol{a} = \boldsymbol{H}\boldsymbol{c}$) that
violate local voltage-power correlations. However, Mahalanobis distance evaluates single snapshots,
making it blind to subtle, gradual ramp drifts that remain within instantaneous confidence ellipsoids.

\section{CUSUM Sequential Change Detection}
Page's Cumulative Sum (CUSUM) test \cite{page1954} integrates deviations across time. The two-sided
recursive formulation maintains positive and negative accumulators \cite{montgomery2019,huang2016}:
\begin{equation}
S_n^+ = \max(0, S_{n-1}^+ + (y_n - \nu_0) - k), \quad S_n^- = \max(0, S_{n-1}^- - (y_n - \nu_0) - k).
\end{equation}
Zero-clamping prevents negative windup during normal operation, allowing instantaneous reaction upon
attack onset. Setting allowance $k = \delta/2$ suppresses Gaussian noise, while threshold $h$
governs the operational trade-off between false-alarm rate ($\text{ARL}_0$) and detection delay ($\text{ARL}_1$).
CUSUM excels at integrating slow ramps, but exhibits multi-frame lag on abrupt shifts.

\section{FPGA-Based Embedded System Implementation}
Substation automation requires line-rate telemetry inspection. Under IEEE~C37.118, PMUs report every
20~ms ($50~\text{Hz}$). General-purpose CPUs and GPUs introduce non-deterministic operating system
jitter and excessive power ($> 100~\text{W}$). In contrast, FPGAs (such as the AMD/Xilinx Artix-7
XC7A200T \cite{artix7ds}) provide deterministic sub-microsecond latency, low power ($< 2~\text{W}$),
and dedicated DSP48E1 arithmetic slices that execute pipelined matrix-vector operations with zero logic bloat.

\section{Fixed-Point Arithmetic and Q-Format Representation}
Floating-point units incur heavy resource and latency penalties. This design adopts a signed
\textbf{Q1.15} fixed-point format (1 sign bit, 15 fractional bits, dynamic range $[-1.0, +0.99997]$),
yielding a theoretical SQNR of $92.07~\text{dB}$. To eliminate intermediate overflow during 4-element
dot products, internal accumulators incorporate 2 guard bits (\textbf{34-bit Q4.30}). Convergent
$+2^{14}$ half-up rounding eliminates DC bias, and saturation clamping prevents two's-complement overflow.

\section{Review of Related Empirical Studies}
Recent literature highlights distinct research directions. Huang \textit{et al.} \cite{huang2016}
and Kumar \textit{et al.} \cite{kumar2022} evaluated CUSUM for gradual drifts, but evaluated
channels independently in software. Kurt \textit{et al.} \cite{kurt2020} demonstrated a 94.2\%
detection rate using decentralized Mahalanobis distance on IEEE 30-bus, but noted high miss rates on
slow drifts. Musleh \textit{et al.} \cite{musleh2023} deployed LSTM autoencoders, achieving 91.4\%
accuracy at the cost of 52~ms inference latency on high-end GPUs. Cui \textit{et al.} \cite{cui2024}
implemented statistical PCA on a multicore DSP, achieving 12~ms latency. Zhang \textit{et al.} \cite{zhang2024}
demonstrated FPGA synchrophasor processing, but restricted evaluation to random noise dropouts.

\section{Comparative Analysis of Existing Approaches}
Table~\ref{tab:priorwork} compares representative published works with the proposed architecture.

\begin{table}[H]
\centering
\caption{Comparison of related works in smart grid FDI anomaly detection.}
\label{tab:priorwork}
\small
\renewcommand{\arraystretch}{1.25}
\begin{tabular}{|p{2.6cm}|p{2.4cm}|p{2.2cm}|p{3.2cm}|p{3.6cm}|}
\hline
\textbf{Author / Year} & \textbf{Detection Method} & \textbf{Platform} & \textbf{Reported Performance} & \textbf{Primary Limitation} \\
\hline
Liu, Ning \& Reiter (2009) \cite{liu2009} & Classical Bad-Data Detection (WLS / $\chi^2$) & Software (MATLAB) & Evaded by $\boldsymbol{a}=\boldsymbol{H}\boldsymbol{c}$ (0\% stealth detection) & Completely blind to coordinated stealth FDI; residual $\boldsymbol{r}$ is invariant. \\
\hline
Huang \textit{et al.} (2016) \cite{huang2016} & Adaptive Cumulative Sum (CUSUM) & Software (MATLAB) & $\text{ARL}_0 = 200$, $\text{ARL}_1 \approx 15$ steps on slow drift & Evaluates channels independently; misses spatial cross-correlation; software latency. \\
\hline
Kurt, Y\i lmaz \& Wang (2020) \cite{kurt2020} & Decentralized Mahalanobis Distance & Software (Python, IEEE 30-bus) & 94.2\% detection rate on abrupt FDI; sub-second & Evaluates single snapshots without memory; blind to slow ramp drifts below threshold. \\
\hline
Musleh, Chen \& Dong (2023) \cite{musleh2023} & Deep Learning (LSTM Autoencoder) & Software / GPU (PyTorch, RTX GPU) & 91.4\% accuracy, F1 = 0.897 on IEEE 14-bus & Latency $> 52$~ms; high power ($> 150$~W); black box lacking statistical guarantees. \\
\hline
Cui, Han \& Zou (2024) \cite{cui2024} & Statistical PCA + Chi-Square & Embedded DSP (TI TMS320C6678) & 88.2\% detection rate; 12~ms latency & High power ($> 10$~W); floating-point bottleneck; vulnerable to stealth subspace FDI. \\
\hline
\textbf{This Work (Group 9, 2026)} & \textbf{Dual-Layer Mahalanobis + Two-Sided CUSUM} & \textbf{AMD/Xilinx Artix-7 FPGA (XC7A200T)} & \textbf{84.67\% accuracy, 96.36\% precision, 71.57\% recall, F1 = 0.8214, 360~ns latency} & Fixed-point precision bounded; unobservable replay attacks require active physical watermarking. \\
\hline
\end{tabular}
\end{table}

\section{Identified Research Gap}
Existing literature reveals a critical gap: theoretical algorithms in software ignore strict
substation line-rate execution constraints, while existing hardware prototypes often rely on
simplistic univariate thresholding that remains vulnerable to coordinated stealth attacks ($\boldsymbol{a} = \boldsymbol{H}\boldsymbol{c}$).
Furthermore, prior studies frequently neglect fixed-point quantization dynamics and evaluate only a
single attack profile. A clear requirement exists for a hardware-efficient, dual-layer statistical
anomaly engine coupling spatial Mahalanobis distance with temporal two-sided CUSUM on FPGA, verified
across multiple adversarial attack regimes with rigorous fixed-point safeguards.

\section{Chapter Summary}
This chapter reviewed smart-grid architecture, the mathematical stealth FDI condition, and the
failure of conventional bad-data detection. The complementary properties of spatial Mahalanobis
distance and temporal two-sided CUSUM were analyzed, and FPGA fixed-point implementation considerations
were detailed. A comparative analysis and research gap established the foundation for the proposed
dual-layer hardware architecture, formulated mathematically in Chapter Three.


% =====================================================================
% CHAPTER THREE
% =====================================================================
\chapter{Mathematical Formulation and Fixed-Point Design}

\section{System Model and Mathematical Notation}
Power system monitoring relies on continuous acquisition of synchronized telemetry
representing the physical operating state of the electrical transmission network.
To establish mathematical rigor across both algorithmic formulation and digital
hardware implementation, Table~\ref{tab:notation} defines the principal symbols,
their physical meanings, coordinate formats, and operational units utilized
throughout this report.

\begin{table}[H]
\centering
\caption{Notation and mathematical symbols used throughout the report.}
\label{tab:notation}
\small
\renewcommand{\arraystretch}{1.25}
\begin{tabular}{|p{2.2cm}|p{7.2cm}|p{2.6cm}|p{2.2cm}|}
\hline
\textbf{Symbol} & \textbf{Mathematical Definition / Physical Meaning} & \textbf{Data Format} & \textbf{Physical Unit} \\
\hline
$y_t$ & Scalar monitored measurement or observation at time $t$ & Signed Q1.15 & p.u. (normalized) \\
\hline
$D_t^2$ & Squared Mahalanobis distance statistic at time $t$ & Signed Q1.15 & Dimensionless \\
\hline
$S_t$ & Cumulative Sum (CUSUM) statistic at time $t$ & Signed Q1.15 & Dimensionless \\
\hline
$h$ & CUSUM decision alarm threshold & Signed Q1.15 & Dimensionless \\
\hline
$k$ & CUSUM reference allowance / slack parameter ($k = \delta/2$) & Signed Q1.15 & Dimensionless \\
\hline
$\boldsymbol{\Sigma}$ & Feature covariance matrix ($d \times d$) & Double / Q1.15 & $\text{p.u.}^2$ \\
\hline
$\boldsymbol{\Sigma}^{-1}, \boldsymbol{W}$ & Inverse covariance (precision) matrix & Signed Q1.15 & $\text{p.u.}^{-2}$ \\
\hline
$\boldsymbol{r}_t$ & Power system state-estimation measurement residual & Real vector & Physical / p.u. \\
\hline
$\alpha$ & False-alarm probability / significance level & Scalar & Probability \\
\hline
$\boldsymbol{z}_t$ & Raw power-system PMU measurement vector ($m \times 1$) & Vector & kV, A, MW, MVAR \\
\hline
$\boldsymbol{x}_t$ & True power-system complex state vector ($n \times 1$) & Vector & Radian, p.u. \\
\hline
$\hat{\boldsymbol{x}}_t$ & Weighted Least-Squares estimated system state & Vector & Radian, p.u. \\
\hline
$\boldsymbol{H}$ & Power system measurement Jacobian / topology matrix & Matrix & Model-dependent \\
\hline
$\boldsymbol{a}_t$ & False Data Injection perturbation vector ($\boldsymbol{a} = \boldsymbol{H}\boldsymbol{c}$) & Vector & Physical / p.u. \\
\hline
$\boldsymbol{c}$ & Attacker false state bias vector & Vector & Radian, p.u. \\
\hline
$\boldsymbol{\mu}$ & Nominal reference feature mean vector ($d \times 1$) & Signed Q1.15 & p.u. (normalized) \\
\hline
$\boldsymbol{q}_t$ & Preprocessed, mean-centered 4-channel feature vector & Signed Q1.15 & p.u. (normalized) \\
\hline
$\tau^2$ & Spatial Mahalanobis decision threshold ($\chi^2_{d, 1-\alpha}$) & Signed Q1.15 & Dimensionless \\
\hline
$S_t^+, S_t^-$ & Two-sided positive and negative recursive CUSUM accumulators & Signed Q1.15 & Dimensionless \\
\hline
$\text{ARL}_0$ & In-control Average Run Length (mean time between false alarms)& Integer & Sample frames \\
\hline
$\text{ARL}_1$ & Out-of-control Average Run Length (mean detection delay) & Integer & Sample frames \\
\hline
\end{tabular}
\end{table}

\section{False Data Injection Model and the Stealth Condition}
The transmission grid telemetry acquisition process is formulated under the linear DC
power-flow approximation or linearized AC operating point:
\begin{equation}
\boldsymbol{z} = \boldsymbol{H}\boldsymbol{x} + \boldsymbol{e},
\label{eq:ch3_dc_model}
\end{equation}
where $\boldsymbol{z} \in \mathbb{R}^m$ is the vector of $m$ measurements, $\boldsymbol{x} \in \mathbb{R}^n$
is the state vector of $n$ bus voltage angles (with $m > n$ to ensure system observability),
and $\boldsymbol{e} \sim \mathcal{N}(\boldsymbol{0}, \boldsymbol{R})$ represents independent,
zero-mean Gaussian measurement errors with known diagonal covariance matrix
$\boldsymbol{R} = \text{diag}(\sigma_1^2, \dots, \sigma_m^2)$.

An adversary executing a False Data Injection (FDI) attack introduces an additive
perturbation vector $\boldsymbol{a} \in \mathbb{R}^m$ into the telemetry stream:
\begin{equation}
\boldsymbol{z}_a = \boldsymbol{z} + \boldsymbol{a}.
\label{eq:ch3_attack_stream}
\end{equation}
The attack is defined as \textbf{coordinated stealth FDI} if the adversary constructs
$\boldsymbol{a}$ as a linear combination of the columns of the measurement matrix $\boldsymbol{H}$:
\begin{equation}
\boldsymbol{a} = \boldsymbol{H}\boldsymbol{c}, \quad \text{for some non-zero } \boldsymbol{c} \in \mathbb{R}^n,
\label{eq:ch3_stealth_definition}
\end{equation}
where $\boldsymbol{c}$ denotes an arbitrary false state bias vector deliberately injected
into the power system operating state.

\section{Weighted Least-Squares State Estimation and Residual Invariance}
The Weighted Least-Squares (WLS) state estimator computes the state estimate $\hat{\boldsymbol{x}}$
that minimizes the weighted sum of squared residuals:
\begin{equation}
J(\boldsymbol{x}) = [\boldsymbol{z} - \boldsymbol{H}\boldsymbol{x}]^T \boldsymbol{R}^{-1} [\boldsymbol{z} - \boldsymbol{H}\boldsymbol{x}].
\label{eq:ch3_wls_cost}
\end{equation}
Differentiating Equation~\eqref{eq:ch3_wls_cost} with respect to $\boldsymbol{x}$ and equating
to zero yields the normal equations:
\begin{equation}
\boldsymbol{H}^T \boldsymbol{R}^{-1} [\boldsymbol{z} - \boldsymbol{H}\hat{\boldsymbol{x}}] = \boldsymbol{0} \implies \hat{\boldsymbol{x}} = (\boldsymbol{H}^T \boldsymbol{R}^{-1} \boldsymbol{H})^{-1} \boldsymbol{H}^T \boldsymbol{R}^{-1} \boldsymbol{z}.
\label{eq:ch3_normal_eqs}
\end{equation}
Let $\boldsymbol{G} = \boldsymbol{H}^T \boldsymbol{R}^{-1} \boldsymbol{H}$ denote the symmetric,
positive-definite power-system gain matrix. The state estimation projection is $\hat{\boldsymbol{x}} = \boldsymbol{G}^{-1} \boldsymbol{H}^T \boldsymbol{R}^{-1} \boldsymbol{z}$.

Under nominal clean operation, the measurement residual vector is:
\begin{equation}
\boldsymbol{r} = \boldsymbol{z} - \boldsymbol{H}\hat{\boldsymbol{x}} = [\boldsymbol{I} - \boldsymbol{H}\boldsymbol{G}^{-1}\boldsymbol{H}^T\boldsymbol{R}^{-1}]\boldsymbol{z} = (\boldsymbol{I} - \boldsymbol{K})\boldsymbol{z},
\label{eq:ch3_clean_residual}
\end{equation}
where $\boldsymbol{K} = \boldsymbol{H}\boldsymbol{G}^{-1}\boldsymbol{H}^T\boldsymbol{R}^{-1}$ is the
orthogonal projection (hat) matrix.

\begin{quote}
\textbf{Theorem 1 (Residual Invariance under Stealth FDI):}
\textit{Let $\boldsymbol{z}_a = \boldsymbol{z} + \boldsymbol{H}\boldsymbol{c}$ represent an attacked
measurement vector. The post-attack state estimate is $\hat{\boldsymbol{x}}_a = \hat{\boldsymbol{x}} + \boldsymbol{c}$,
and the resulting post-attack residual $\boldsymbol{r}_a$ is strictly identical to the pre-attack residual $\boldsymbol{r}$.}
\end{quote}

\noindent\textbf{Proof:}
Applying the linear WLS estimator to the corrupted vector $\boldsymbol{z}_a$:
\begin{align}
\hat{\boldsymbol{x}}_a &= \boldsymbol{G}^{-1} \boldsymbol{H}^T \boldsymbol{R}^{-1} \boldsymbol{z}_a \nonumber \\
&= \boldsymbol{G}^{-1} \boldsymbol{H}^T \boldsymbol{R}^{-1} (\boldsymbol{z} + \boldsymbol{H}\boldsymbol{c}) \nonumber \\
&= \boldsymbol{G}^{-1} \boldsymbol{H}^T \boldsymbol{R}^{-1} \boldsymbol{z} + \boldsymbol{G}^{-1} (\boldsymbol{H}^T \boldsymbol{R}^{-1} \boldsymbol{H})\boldsymbol{c} \nonumber \\
&= \hat{\boldsymbol{x}} + \boldsymbol{G}^{-1}\boldsymbol{G}\boldsymbol{c} \nonumber \\
&= \hat{\boldsymbol{x}} + \boldsymbol{c}.
\end{align}
Evaluating the post-attack measurement residual:
\begin{align}
\boldsymbol{r}_a &= \boldsymbol{z}_a - \boldsymbol{H}\hat{\boldsymbol{x}}_a \nonumber \\
&= (\boldsymbol{z} + \boldsymbol{H}\boldsymbol{c}) - \boldsymbol{H}(\hat{\boldsymbol{x}} + \boldsymbol{c}) \nonumber \\
&= \boldsymbol{z} + \boldsymbol{H}\boldsymbol{c} - \boldsymbol{H}\hat{\boldsymbol{x}} - \boldsymbol{H}\boldsymbol{c} \nonumber \\
&= \boldsymbol{z} - \boldsymbol{H}\hat{\boldsymbol{x}} \nonumber \\
&= \boldsymbol{r}.
\end{align}
Evaluating the global chi-square detection statistic:
\begin{equation}
J(\hat{\boldsymbol{x}}_a) = \boldsymbol{r}_a^T \boldsymbol{R}^{-1} \boldsymbol{r}_a = \boldsymbol{r}^T \boldsymbol{R}^{-1} \boldsymbol{r} = J(\hat{\boldsymbol{x}}).
\end{equation}
Thus, $J(\hat{\boldsymbol{x}}_a)$ is mathematically identical to $J(\hat{\boldsymbol{x}})$. The
chi-square test fails to detect the attack, completing the proof. \hfill $\blacksquare$

\section{Spatial Mahalanobis-Distance Detector Formulation}
Because the state-estimation residual is invariant under $\boldsymbol{a} = \boldsymbol{H}\boldsymbol{c}$,
detection must exploit the multivariate statistical relationships governing local physical
quantities. Let $\boldsymbol{q}_t \in \mathbb{R}^d$ ($d=4$) represent the localized synchrophasor
feature vector capturing terminal bus voltages and active/reactive branch flows ($V_1, V_2, P_{12}, Q_{12}$)
at sampling instant $t$.

During a clean, unattacked baseline calibration period $t \in [1, T_{\text{cal}}]$, the nominal
feature mean vector $\boldsymbol{\mu} \in \mathbb{R}^d$ and sample covariance matrix
$\boldsymbol{\Sigma} \in \mathbb{R}^{d \times d}$ are computed:
\begin{align}
\boldsymbol{\mu} &= \frac{1}{T_{\text{cal}}} \sum_{t=1}^{T_{\text{cal}}} \boldsymbol{q}_t, \label{eq:ch3_mean_est} \\
\boldsymbol{\Sigma} &= \frac{1}{T_{\text{cal}} - 1} \sum_{t=1}^{T_{\text{cal}}} (\boldsymbol{q}_t - \boldsymbol{\mu})(\boldsymbol{q}_t - \boldsymbol{\mu})^T. \label{eq:ch3_cov_est}
\end{align}
To eliminate the requirement for online matrix inversion inside the hardware datapath, the
precision matrix $\boldsymbol{W} = \boldsymbol{\Sigma}^{-1}$ is precomputed offline and loaded
into static FPGA configuration registers.

For each incoming observation $\boldsymbol{q}_t$, the preprocessed deviation vector is
$\Delta\boldsymbol{q}_t = \boldsymbol{q}_t - \boldsymbol{\mu}$. The instantaneous squared
Mahalanobis distance statistic is:
\begin{equation}
D_t^2 = \Delta\boldsymbol{q}_t^T \boldsymbol{W} \Delta\boldsymbol{q}_t = \sum_{i=1}^d \sum_{j=1}^d \Delta q_{t,i} W_{ij} \Delta q_{t,j}.
\label{eq:ch3_mahal_stat}
\end{equation}
Under the null hypothesis of uncorrupted, stationary Gaussian operation ($H_0: \boldsymbol{q}_t \sim \mathcal{N}(\boldsymbol{\mu}, \boldsymbol{\Sigma})$),
the quadratic statistic $D_t^2$ is distributed as a central chi-square random variable with
$d = 4$ degrees of freedom:
\begin{equation}
D_t^2 \sim \chi^2_d.
\end{equation}
The Layer 1 spatial decision rule asserts an alarm flag $A_{M,t}$ whenever $D_t^2$ exceeds
the calibrated threshold $\tau^2$:
\begin{equation}
A_{M,t} = \begin{cases}
1, & \text{if } D_t^2 > \tau^2, \\
0, & \text{if } D_t^2 \le \tau^2,
\end{cases}
\quad \text{where } \Pr(D_t^2 > \tau^2 \mid H_0) = \alpha.
\label{eq:ch3_layer1_decision}
\end{equation}

\section{Temporal Two-Sided CUSUM Sequential Detector Formulation}
While Layer 1 provides single-cycle detection for abrupt and scaling attacks, slow ramp
attacks with small injection rates ($\boldsymbol{a}_t = \beta t \boldsymbol{H}\boldsymbol{c}$)
produce instantaneous deviations that remain below $\tau^2$ during initial phases.
Layer 2 addresses this vulnerability by deploying a recursive, two-sided Cumulative Sum
(CUSUM) filter to integrate subtle persistent shifts over time.

Let $y_{c,t}$ denote the normalized scalar telemetry signal for channel $c \in \{1, \dots, d\}$
at time $t$, with nominal reference mean $\nu_{0,c}$. The two-sided recursive CUSUM
formulation tracks positive shifts ($S_{c,t}^+$) and negative shifts ($S_{c,t}^-$):
\begin{align}
S_{c,t}^+ &= \max\left(0, \; S_{c,t-1}^+ + (y_{c,t} - \nu_{0,c}) - k_c\right), \label{eq:ch3_cusum_pos} \\
S_{c,t}^- &= \max\left(0, \; S_{c,t-1}^- - (y_{c,t} - \nu_{0,c}) - k_c\right), \label{eq:ch3_cusum_neg}
\end{align}
with initial boundary conditions $S_{c,0}^+ = S_{c,0}^- = 0$.

The CUSUM parameters are mathematically derived from the target shift magnitude $\delta_c$:
\begin{equation}
k_c = \frac{\delta_c}{2},
\end{equation}
which establishes the drift allowance deadband. The alarm threshold $h_c$ is determined
to balance the in-control run length $\text{ARL}_0$ against out-of-control detection delay $\text{ARL}_1$.

The Layer 2 temporal decision rule asserts an alarm flag $A_{C,c,t}$ for channel $c$ if
either accumulator crosses threshold $h_c$:
\begin{equation}
A_{C,c,t} = \begin{cases}
1, & \text{if } S_{c,t}^+ > h_c \quad \text{or} \quad S_{c,t}^- > h_c, \\
0, & \text{otherwise}.
\end{cases}
\label{eq:ch3_layer2_decision}
\end{equation}

\section{Combined Dual-Layer Decision Rule}
The global anomaly decision flag $A_t \in \{0, 1\}$ at time step $t$ is formulated as the
logical disjunction (OR) of the spatial Mahalanobis alert and all multi-channel temporal
CUSUM alerts:
\begin{equation}
A_t = A_{M,t} \lor \left( \bigvee_{c=1}^d A_{C,c,t} \right).
\label{eq:ch3_global_decision}
\end{equation}
This synergistic formulation guarantees that an alarm is declared immediately if an abrupt
structural perturbation violates spatial cross-channel covariance ($A_{M,t} = 1$) or if a
subtle gradual drift accumulates sufficient temporal evidence to trip CUSUM ($A_{C,c,t} = 1$).

\section{Feature Normalization and Mean-Centering}
Power system telemetry variables encompass vastly different physical units and magnitudes:
voltage magnitudes hover near $1.0~\text{p.u.}$ ($13.8\text{--}138~\text{kV}$), while branch
active power flows can span hundreds of megawatts. In digital fixed-point hardware, raw
unscaled numbers would cause catastrophic dynamic range underflow in fractional channels
and severe overflow in high-power channels.

To establish uniform numerical conditioning, each raw measurement $z_c$ is linearly mapped
into a normalized fixed-point range using nominal baseline operating points:
\begin{equation}
q_c = \frac{z_c - \bar{z}_c}{\sigma_{z,c}} \times \text{ScaleFactor},
\end{equation}
ensuring that all preprocessed features $\Delta q_c = q_c - \mu_c$ conform strictly to the
$[-1.0, +0.999969]$ boundaries of the signed Q1.15 numerical format.

\section{Fixed-Point Signed Q1.15 Arithmetic and Quantization Error Analysis}
Digital signal processing algorithms targeting FPGA fabric face a fundamental trade-off
between numerical precision and silicon resource utilization. Floating-point arithmetic
incurs high DSP slice consumption and multi-cycle pipeline latency. This project implements
a uniform signed \textbf{Q1.15} fixed-point representation for all inputs, calibration
matrices, and arithmetic datapaths.

A 16-bit signed two's-complement word in Q1.15 format consists of 1 sign bit and 15 fractional
bits. Its decimal value is:
\begin{equation}
X = -b_{15} + \sum_{i=0}^{14} b_i 2^{i - 15}, \quad b_i \in \{0, 1\}.
\end{equation}
The quantization step size is $\Delta = 2^{-15} \approx 3.051758 \times 10^{-5}$. Assuming
round-off quantization noise $e_q$ is uniformly distributed over $[-\Delta/2, +\Delta/2]$,
the noise variance is:
\begin{equation}
\sigma_q^2 = \frac{\Delta^2}{12} = \frac{2^{-30}}{12} \approx 7.761 \times 10^{-11}.
\end{equation}
The Signal-to-Quantization-Noise Ratio (SQNR) for a full-scale sinusoidal input is:
\begin{equation}
\text{SQNR} = 6.0206 \times 15 + 1.7609 \approx 92.07~\text{dB}.
\end{equation}
This theoretical noise floor is over $25~\text{dB}$ below the typical $0.1\%$ total vector
error (TVE) threshold mandated by the IEEE~C37.118 synchrophasor standard, proving that Q1.15
fixed-point representation introduces negligible distortion into the anomaly detection process.

\section{Accumulator Word Length, Guard Bits, and Saturation Safeguards}
When evaluating the matrix-vector multiplication $\boldsymbol{y} = \boldsymbol{W} \Delta\boldsymbol{q}$
and the subsequent quadratic dot product $D^2 = \Delta\boldsymbol{q}^T \boldsymbol{y}$, two 16-bit
Q1.15 operands are multiplied, generating an exact 32-bit product:
\begin{equation}
\text{Q1.15} \times \text{Q1.15} \implies \text{Q2.30 (32 bits)}.
\end{equation}
The dot product computes the sum of $d = 4$ products. In the worst-case scenario where all
terms attain full-scale values simultaneously, intermediate summation can cause bit overflow.
To mathematically eliminate overflow, the accumulator word length must include guard bits:
\begin{equation}
N_{\text{guard}} = \lceil \log_2(d) \rceil = \lceil \log_2(4) \rceil = 2~\text{bits}.
\end{equation}
Thus, the internal accumulation register is sized to \textbf{34 bits (Q4.30)}, providing
dynamic range $[-4.0, +3.9999999]$ with 30 fractional bits.

Following accumulation, the 34-bit value is scaled back to 16 bits using convergent half-up
rounding by adding $+2^{14} = \text{0x0000\_4000}$ before truncating bits $[30:15]$. To protect
against two's-complement arithmetic overflow (e.g., $-1.0 \times -1.0 = +1.0$, which exceeds
the maximum positive Q1.15 representable value of $+0.999969$), dedicated saturation logic
clamps any overflow to \texttt{0x7FFF} (+0.999969) and any underflow to \texttt{0x8000} (-1.0).

\section{Project Calibration Parameters and Fixed-Point Wordlengths}
Table~\ref{tab:q15_params} summarizes the complete fixed-point parameter configuration,
word lengths, quantization boundaries, and architectural assignments implemented in the
RTL design.

\begin{table}[H]
\centering
\caption{Fixed-point Q1.15 design parameters and quantization boundaries.}
\label{tab:q15_params}
\small
\renewcommand{\arraystretch}{1.25}
\begin{tabular}{|p{2.6cm}|p{4.2cm}|p{2.0cm}|p{1.4cm}|p{2.4cm}|p{2.4cm}|}
\hline
\textbf{Parameter} & \textbf{Physical Description} & \textbf{Format} & \textbf{Bits} & \textbf{Scaling ($2^{-15}$)} & \textbf{Dynamic Range} \\
\hline
$z_c$ & PMU input telemetry vector & Signed Q1.15 & 16 & $3.0518 \times 10^{-5}$ & $[-1.0, +0.99997]$ \\
\hline
$\mu_c$ & Historical baseline mean vector & Signed Q1.15 & 16 & $3.0518 \times 10^{-5}$ & $[-1.0, +0.99997]$ \\
\hline
$r_c = z_c - \mu_c$ & Preprocessed deviation residual & Signed Q1.15 & 16 & $3.0518 \times 10^{-5}$ & $[-1.0, +0.99997]$ \\
\hline
$W_{ij}$ & Inverse covariance matrix entry & Signed Q1.15 & 16 & $3.0518 \times 10^{-5}$ & $[-1.0, +0.99997]$ \\
\hline
$P_{\text{mult}}$ & DSP multiplier product & Signed Q2.30 & 32 & $9.3132 \times 10^{-10}$ & $[-2.0, +1.99999]$ \\
\hline
$A_{\text{dot}}$ & Dot-product accumulator & Signed Q4.30 & 34 & $9.3132 \times 10^{-10}$ & $[-8.0, +7.99999]$ \\
\hline
$D^2$ & Squared Mahalanobis distance & Signed Q1.15 & 16 & $3.0518 \times 10^{-5}$ & $[0.0, +0.99997]$ \\
\hline
$\tau^2$ & Spatial decision threshold & Signed Q1.15 & 16 & $1673 \times 2^{-15}$ & $0.05106~\text{p.u.}$ \\
\hline
$k$ & CUSUM reference slack allowance & Signed Q1.15 & 16 & $5461 \times 2^{-15}$ & $0.16667~\text{p.u.}$ \\
\hline
$h$ & CUSUM decision alarm threshold & Signed Q1.15 & 16 & $20753 \times 2^{-15}$ & $0.63333~\text{p.u.}$ \\
\hline
\end{tabular}
\end{table}

\section{Hardware Latency and Throughput Closed-Form Mathematical Model}
To provide analytical performance guarantees prior to hardware synthesis, closed-form
equations model execution latency, processing throughput, and power dissipation.

Let $f_{\text{clk}}$ denote the system clock frequency and $T_{\text{clk}} = 1/f_{\text{clk}}$
denote the clock period. The total execution latency $T_{\text{latency}}$ to process a single
4-channel measurement frame through the FSM controller and pipelined datapath is:
\begin{equation}
T_{\text{latency}} = \frac{N_{\text{cycles}}}{f_{\text{clk}}},
\label{eq:ch3_latency_model}
\end{equation}
where $N_{\text{cycles}}$ is the cycle count derived from the state machine schedule.
For the 4-channel engine operating at $f_{\text{clk}} = 100~\text{MHz}$ ($T_{\text{clk}} = 10~\text{ns}$)
with $N_{\text{cycles}} = 36$~cycles:
\begin{equation}
T_{\text{latency}} = 36 \times 10~\text{ns} = 360~\text{ns} = 0.36~\mu\text{s}.
\end{equation}

The theoretical maximum processing throughput in measurement frames per second is:
\begin{equation}
\text{Throughput} = \frac{f_{\text{clk}}}{N_{\text{cycles}}} = \frac{100 \times 10^6}{36} \approx 2,777,777~\text{frames/s} \approx 2.78~\text{MSamples/s}.
\label{eq:ch3_throughput_model}
\end{equation}

The maximum attainable operating clock frequency $f_{\max}$ is bounded by the critical
path delay $T_{\text{crit}}$ and Worst Negative Slack (WNS):
\begin{equation}
f_{\max} = \frac{1}{T_{\text{clk}} - \text{WNS}}.
\label{eq:ch3_fmax_model}
\end{equation}

The dynamic energy consumed per anomaly detection operation $E_{\text{op}}$ is given by:
\begin{equation}
E_{\text{op}} = P_{\text{dyn}} \times T_{\text{latency}} = P_{\text{dyn}} \times \left( \frac{N_{\text{cycles}}}{f_{\text{clk}}} \right),
\label{eq:ch3_energy_model}
\end{equation}
where $P_{\text{dyn}}$ is the dynamic power reported by Vivado power analysis. These analytical
models are verified experimentally against Vivado synthesis reports in Chapter Five.


% =====================================================================
% CHAPTER FOUR
% =====================================================================
\chapter{Methodology and Hardware Implementation}

\section{Overall Methodology and Repeatable 10-Step Design Flow}
To ensure technical rigor, repeatability, and complete traceability between high-level
mathematics and physical digital hardware, the engineering methodology follows a structured
10-step design flow:
\begin{enumerate}
    \item \textbf{Step 1: System Specification:} Define functional requirements, input telemetry
    formats (signed Q1.15), target reporting rates ($F_s = 50~\text{Hz}$ under IEEE~C37.118),
    and threat models (stealth, ramp, scaling, and replay FDI).
    \item \textbf{Step 2: Floating-Point Algorithmic Model:} Implement a software prototype in
    Python to formulate the mathematical state estimator, Mahalanobis distance, and two-sided
    CUSUM algorithms.
    \item \textbf{Step 3: Bit-Accurate Fixed-Point Golden Model:} Translate the floating-point
    model into an exact, integer-scaled Q1.15 bit-accurate golden model in Python. Every hardware
    output is verified against this reference.
    \item \textbf{Step 4: Microarchitecture Design:} Develop block diagrams, state transition
    graphs, memory maps, and datapath schedules. Choose between parallel and time-multiplexed
    hardware topologies.
    \item \textbf{Step 5: Synthesizable RTL Description:} Code the design in SystemVerilog
    (IEEE 1800-2017) using fully parameterizable word lengths and strictly synthesizable constructs.
    \item \textbf{Step 6: Functional \& Assertion-Based Verification:} Construct self-checking
    SystemVerilog testbenches with known-answer test (KAT) vectors and automated mismatch detection.
    \item \textbf{Step 7: Logic Synthesis and Technology Mapping:} Synthesize the RTL in AMD/Xilinx
    Vivado in accordance with Vivado synthesis and constraint guidelines \cite{ug901,ug903,ug907} targeting the Artix-7 XC7A200T FPGA; extract resource utilization (LUTs, FFs, DSPs, BRAM).
    \item \textbf{Step 8: Timing and Power Analysis:} Evaluate static timing slack (WNS/WHS),
    maximum clock frequency ($f_{\max}$), and dynamic power consumption via Vivado report tools.
    \item \textbf{Step 9: Automated Telemetry \& CSV Data Export:} Execute automated regression
    scripts that write hardware outputs to CSV format for statistical performance analysis.
    \item \textbf{Step 10: Reproducibility Package Assembly:} Archive all RTL, testbenches, Python
    scripts, synthesis Tcl files, and raw simulation datasets into an open, documented repository.
\end{enumerate}

\section{Python Telemetry Generation and IEEE 14-Bus Test System}
To evaluate the anomaly engine under realistic power grid operating conditions, telemetry
streams are derived from the standard IEEE 14-bus transmission network benchmark \cite{matpowercase14}.
The IEEE 14-bus system represents a portion of the Midwestern American electric power system,
comprising 5 synchronous generators, 14 transmission buses, and 20 interconnecting lines.

The Python data generator executes a steady-state Newton-Raphson power-flow solution using
MATPOWER reference formulations \cite{zimmerman2011,matpowercase14}. A four-dimensional synchrophasor measurement configuration
is extracted:
\begin{equation}
\boldsymbol{z}_t = [V_{1,t}, \; V_{2,t}, \; P_{12,t}, \; Q_{12,t}]^T,
\end{equation}
where $V_1, V_2$ represent voltage magnitudes at Bus 1 (slack bus) and Bus 2 (generator bus),
and $P_{12}, Q_{12}$ denote the active and reactive power flows traversing the high-voltage
transmission branch connecting Bus 1 and Bus 2. Stationary Gaussian measurement noise is
superimposed onto all channels ($\sigma = 0.005\text{--}0.020~\text{p.u.}$) to emulate physical
instrument transformer and PMU analog-to-digital conversion noise.

\section{Cyber-Attack Injection Simulation}
Across the 1,200 continuous frames generated at $F_s = 50~\text{Hz}$ ($20~\text{ms}$ interval,
$24.0~\text{s}$ total sequence), four distinct cyber-attack modalities are injected:
\begin{itemize}
    \item \textbf{Scenario 1: Coordinated Stealth FDI ($t = 200\text{--}349$, 150 frames):}
    The adversary computes the measurement Jacobian $\boldsymbol{H}$ and injects a covert
    perturbation $\boldsymbol{a} = \boldsymbol{H}\boldsymbol{c}$, where $\boldsymbol{c}$ applies
    a false $+0.035~\text{rad}$ phase angle bias. Active power flow $P_{12}$ shifts by $+0.045~\text{p.u.}$
    while $V_2$ shifts by $-0.015~\text{p.u.}$, perfectly evading classical $\chi^2$ bad data detection.
    \item \textbf{Scenario 2: Gradual Ramp / Drift Attack ($t = 450\text{--}649$, 200 frames):}
    A subtle linear ramp is introduced at a rate of $+0.0004~\text{p.u./frame}$ ($+0.08~\text{p.u.}$ total).
    The gradual slope bypasses instantaneous threshold checkers, testing the CUSUM integration capability.
    \item \textbf{Scenario 3: Multi-Channel False Scaling Attack ($t = 750\text{--}869$, 120 frames):}
    Telemetry is scaled by a multiplicative factor of $+5.0\%$ ($(1 + \alpha)\boldsymbol{z}_t$, $\alpha = 0.05$),
    amplifying both steady-state loading and high-frequency noise.
    \item \textbf{Scenario 4: Historical Replay Attack ($t = 950\text{--}1070$, 121 frames):}
    Telemetry recorded during clean historical frames ($t = 50\text{--}170$) is replayed during an
    actual physical grid load excursion, masking the true operational state from operators.
\end{itemize}

\section{Calibration and Evaluation Dataset Split}
To avoid methodological bias and data leakage, the 1,200-frame time series is partitioned into
strict non-overlapping operational windows:
\begin{itemize}
    \item \textbf{Calibration Window (Frames 0..199, 200 frames / 4.0 s):} Exclusively clean,
    stationary baseline telemetry. Used to compute reference mean vector $\boldsymbol{\mu}$,
    sample covariance $\boldsymbol{\Sigma}$, and inverse precision matrix $\boldsymbol{W} = \boldsymbol{\Sigma}^{-1}$.
    Zero attack contamination is permitted.
    \item \textbf{Test Evaluation Split (Frames 200..1199, 1,000 frames / 20.0 s):} Contains all
    four attack injection windows interspersed with unattacked recovery periods. All statistical
    detection metrics (Accuracy, Precision, Recall, F1, ROC-AUC) are evaluated on this test split.
\end{itemize}

\section{Hardware Architecture and Module Partitioning}
Figure~\ref{fig:hardware_choose} illustrates the structural partitioning between software
offline calibration and dedicated FPGA hardware execution.

\reportfigure[0.85\textwidth]{hardware_choose}{Hardware architecture selection and processing partition.}{fig:hardware_choose}

The hardware architecture (\texttt{pmu\_anomaly\_top.sv}) decomposes the anomaly detection engine
into dedicated, interacting functional blocks:
\begin{itemize}
    \item \textbf{Master FSM Controller (\texttt{pmu\_fsm\_ctrl}):} Coordinates synchronous sample
    capture, datapath handshakes, threshold evaluations, and state memory write-backs.
    \item \textbf{Mahalanobis Matrix Engine (\texttt{mahalanobis\_matrix\_engine}):} Computes spatial
    matrix-vector multiplication and quadratic dot product using dedicated DSP48E1 slices.
    \item \textbf{Two-Sided CUSUM Core (\texttt{cusum\_filter\_q15}):} Implements signed positive and
    negative recursive accumulations with zero-clamping logic.
    \item \textbf{Time-Multiplexed State Memory (\texttt{pmu\_state\_bram}):} Stores per-channel CUSUM
    accumulator states and history flags in dual-port Block RAM.
\end{itemize}

Figure~\ref{fig:rtl_schematic} displays the elaborated RTL schematic generated by Vivado
Design Suite, illustrating the inter-module signal connectivity and bus widths.

\reportfigure[0.65\textwidth]{circuit_diagram}{Elaborated top-level RTL schematic of the PMU anomaly engine.}{fig:rtl_schematic}

\section{Circuit-Level Datapath and Arithmetic Safeguards}
Figure~\ref{fig:circuit_diagram} shows the circuit-level datapath schematic, highlighting
the internal arithmetic structures and numerical protection stages.

\reportfigure[0.65\textwidth]{schematic}{Circuit-level schematic of the fixed-point arithmetic datapath.}{fig:circuit_diagram}

The datapath incorporates four critical hardware safeguards:
\begin{enumerate}
    \item \textbf{DSP48E1 Multiplier Utilization:} 16-bit signed operands multiply directly
    into 32-bit Q2.30 intermediate products using hardware DSP slices, preventing fabric logic bloat.
    \item \textbf{34-Bit Accumulator Guard Bits:} An internal 34-bit accumulator (Q4.30) provides
    2 integer guard bits ($32 + \lceil \log_2 4 \rceil = 34$), eliminating intermediate overflow
    during 4-element dot-product summations.
    \item \textbf{Convergent Half-Up Rounding:} Rescaling from 34 bits back to 16 bits adds
    $+2^{14} = \text{0x0000\_4000}$ before truncating bits $[30:15]$, eliminating negative DC bias.
    \item \textbf{Two's-Complement Saturation Logic:} Prevents negative overflow wrapping;
    excursions beyond $+0.999969$ clamp strictly to \texttt{0x7FFF}, and negative excursions clamp to \texttt{0x8000}.
\end{enumerate}

\section{Main Controller Finite-State Machine}
The operational sequencing of the hardware engine is governed by a canonical three-process
Moore Finite State Machine implemented in \texttt{rtl/pmu\_fsm\_ctrl.sv}. Process~1 executes the
sequential state register with asynchronous active-low reset; Process~2 evaluates combinational
next-state transition logic; and Process~3 generates pure Moore output strobes derived
strictly from the current state. Table~\ref{tab:fsm_states} details the complete operational
sequence, matching the actual RTL implementation.

\begin{table}[H]
\centering
\caption{FSM states and operational sequence of the PMU controller.}
\label{tab:fsm_states}
\small
\renewcommand{\arraystretch}{1.25}
\begin{tabular}{|p{3.2cm}|p{3.8cm}|p{7.4cm}|}
\hline
\textbf{FSM State} & \textbf{Architectural Function} & \textbf{Hardware Action Executed in RTL} \\
\hline
\texttt{FSM\_IDLE} & Wait for input sample & Asserts \texttt{data\_ready = 1'b1}; awaits upstream \texttt{data\_valid} strobe from PMU stream. \\
\hline
\texttt{FSM\_CAPTURE} & Latch measurement vector & Asserts \texttt{latch\_sample\_o = 1'b1}; latches 64-bit parallel PMU telemetry vector into input pipeline. \\
\hline
\texttt{FSM\_PREPROCESS} & Mean-centering \& residual & Asserts \texttt{preprocess\_en\_o = 1'b1}; computes signed deviation residual $\boldsymbol{r} = \boldsymbol{z} - \boldsymbol{\mu}$ using calibration registers. \\
\hline
\texttt{FSM\_MAHALANOBIS} & Compute spatial distance & Asserts \texttt{mahal\_start\_o = 1'b1}; triggers 3-cycle matrix engine to compute $\boldsymbol{y} = \boldsymbol{W}\boldsymbol{r}$ and $D^2 = \boldsymbol{r}^T\boldsymbol{y}$; waits for \texttt{mahal\_done}. \\
\hline
\texttt{FSM\_CUSUM} & Update CUSUM filter & Asserts \texttt{cusum\_en\_o = 1'b1}; updates positive ($S^+$) and negative ($S^-$) sequential accumulators with zero-clamping. \\
\hline
\texttt{FSM\_DECISION} & Compare thresholds & Asserts \texttt{decision\_en\_o = 1'b1}; evaluates $D^2 \gtrless \tau^2$ and $S^+, S^- \gtrless h$; asserts global alarm if breached. \\
\hline
\texttt{FSM\_STORE\_RESULT} & Write-back state to BRAM & Asserts \texttt{bram\_we\_o = 1'b1}; writes updated channel CUSUM accumulators and status flags to time-multiplexed BRAM. \\
\hline
\texttt{FSM\_NEXT\_CHANNEL} & Channel advance / loop & Asserts \texttt{chan\_advance\_o = 1'b1}; asserts \texttt{frame\_done\_o = 1'b1} when \texttt{channel\_last} is high; loops to \texttt{CAPTURE} or \texttt{IDLE}. \\
\hline
\end{tabular}
\end{table}

\section{Mahalanobis Matrix Engine Datapath}
The Mahalanobis engine (\texttt{rtl/mahalanobis\_matrix\_engine.sv}) computes the quadratic form
$D^2 = \boldsymbol{r}^T \boldsymbol{W} \boldsymbol{r}$ over a fully pipelined 3-cycle sequence:
\begin{itemize}
    \item \textbf{Cycle 1 (Matrix-Vector Multiplication):} Four parallel row-MAC units multiply
    the 16-bit residual vector $\boldsymbol{r} = [r_1, r_2, r_3, r_4]^T$ by the $4 \times 4$
    precomputed precision matrix $\boldsymbol{W}$:
    \begin{equation}
    y_i = \sum_{j=1}^4 W_{ij} r_j, \quad i \in \{1, 2, 3, 4\}.
    \end{equation}
    Each row multiplication consumes 4 DSP multipliers operating in parallel.
    \item \textbf{Cycle 2 (Quadratic Dot Product):} The resulting intermediate vector $\boldsymbol{y}$
    is multiplied element-wise by the original residual vector $\boldsymbol{r}$:
    \begin{equation}
    D^2 = \sum_{i=1}^4 r_i y_i.
    \end{equation}
    This summation is accumulated in a 34-bit register with 2 guard bits.
    \item \textbf{Cycle 3 (Rounding \& Output Registration):} The 34-bit accumulator value is
    rounded to Q1.15 using convergent half-up rounding, saturated, registered onto the output
    port, and the completion strobe \texttt{mahal\_done} is asserted.
\end{itemize}

\section{Two-Sided CUSUM Sequential RTL Datapath}
The CUSUM core (\texttt{rtl/cusum\_filter\_q15.sv}) executes Equations~\eqref{eq:ch3_cusum_pos} and
\eqref{eq:ch3_cusum_neg} in a single clock cycle. The incoming preprocessed measurement $y_t$
is subtracted from reference mean $\nu_0$ and reference allowance $k$:
\begin{align}
\text{diff\_pos} &= (y_t - \nu_0) - k, \\
\text{diff\_neg} &= -(y_t - \nu_0) - k.
\end{align}
These differences are added to the registered accumulator values $S_{t-1}^+$ and $S_{t-1}^-$.
Dedicated combinational comparators evaluate whether the sum is less than zero; if negative,
the accumulator input is multiplexed to zero, executing instantaneous zero-clamping reset.
Parallel threshold comparators compare $S_t^+$ and $S_t^-$ against $h$, asserting the channel
alarm strobe if either threshold is crossed.

\section{Dual-Port BRAM / LUTRAM State Storage}
The anomaly engine incorporates embedded memory (\texttt{rtl/pmu\_state\_bram.sv}) to maintain CUSUM
accumulator states across successive sample frames. The per-channel state record packs 42 bits:
\begin{itemize}
    \item Positive Accumulator ($S^+$): 16 bits Q1.15.
    \item Negative Accumulator ($S^-$): 16 bits Q1.15.
    \item Channel Alarm Latch: 1 bit.
    \item Anomaly Counter: 8 bits.
    \item Channel Valid Flag: 1 bit.
\end{itemize}
When synthesized targeting dedicated primitives, Vivado infers two RAMB18E1 block RAMs because
the 42-bit packed word exceeds the 36-bit port width of a single RAMB18 block. Alternatively,
the architecture supports parameterization for distributed LUTRAM synthesis, requiring only
11 SLICEM LUTs and zero Block RAMs.

\section{Hardware Timing and Complexity Analysis}
Table~\ref{tab:hardware_timing} presents the verified cycle counts, execution latency,
line-rate throughput, and operational capacity of the hardware anomaly engine operating at
a 100~MHz clock target.

\begin{table}[H]
\centering
\caption{Hardware timing and complexity analysis.}
\label{tab:hardware_timing}
\small
\renewcommand{\arraystretch}{1.25}
\begin{tabular}{|p{4.2cm}|p{3.2cm}|p{2.2cm}|p{1.8cm}|p{3.0cm}|}
\hline
\textbf{Timing / Complexity Parameter} & \textbf{Verification Status} & \textbf{Value} & \textbf{Unit} & \textbf{Operational Meaning} \\
\hline
System Clock Frequency ($f_{\text{clk}}$) & Design Target & 100.0 & MHz & Nominal clock period $T_{\text{clk}} = 10.0~\text{ns}$. \\
\hline
Channel Processing Latency & Verified RTL Cycle Count & 9 & Cycles & Execution time per PMU channel ($90~\text{ns}$). \\
\hline
Full 4-Channel Frame Latency & Verified RTL Cycle Count & 36 & Cycles & Total execution latency ($360~\text{ns} = 0.36~\mu\text{s}$). \\
\hline
Maximum Processing Throughput & Architectural Calculation & 2.78 & MSamples/s & 2,777,777 measurement frames per second. \\
\hline
PMU Standard Period (IEEE C37.118) & Standard Specification & 20.0 & ms & Transmission PMU frame arrival (50~Hz). \\
\hline
Hardware Idle Duty Cycle & Analytical Ratio & 99.9982 & \% & Core remains idle for 19.99964~ms per frame. \\
\hline
Time-Multiplexed PMU Capacity & Architectural Calculation & 55,555 & Streams & Single FPGA core capacity for 50~Hz PMUs. \\
\hline
Worst Negative Slack (WNS) & Estimated Post-Synthesis & Pending & ns & Synthesis clean; pending post-route sign-off. \\
\hline
\end{tabular}
\end{table}

\section{Target FPGA Platform: Artix-7 XC7A200T-FBG676-2}
The AMD/Xilinx Artix-7 XC7A200T device serves as the primary implementation target.
The device characteristics include:
\begin{itemize}
    \item Device Part: \texttt{xc7a200tfbg676-2}.
    \item Logic Slices: 33,650 (134,600 LUT6, 269,200 Registers).
    \item DSP48E1 Multipliers: 740 slices.
    \item Block RAM: 730 RAMB18E1 equivalents (13,140 Kb).
    \item Bonded IOBs: 400 user I/O pins.
\end{itemize}

\section{Verification Strategy and Testbench Hierarchy}
The verification hierarchy comprises two complementary layers:
\begin{enumerate}
    \item \textbf{Unit-Level Testbenches:} Standalone testbenches verify individual modules:
    \texttt{tb\_fixed\_point\_multiplier.sv} verifies rounding and saturation corner cases;
    \texttt{tb\_mahalanobis\_matrix.sv} tests matrix-vector dot products against precomputed
    algebraic vectors; \texttt{tb\_cusum\_filter.sv} tests positive/negative drift clamping.
    \item \textbf{System-Level Testbench (\texttt{tb\_pmu\_anomaly\_top}):} Feeds the entire 1,200-frame
    telemetry sequence from \texttt{pmu\_telemetry.csv} through the top-level RTL entity, logs hardware
    alerts and scores to \texttt{rtl\_output.csv}, and performs automated assertion checks against the
    Python golden model.
\end{enumerate}


% =====================================================================
% CHAPTER FIVE
% =====================================================================
\chapter{Results and Discussion}

\section{Overview of Empirical Results and Discrepancy Reconciliation}
This chapter presents the quantitative empirical results obtained from the Python golden
reference model, self-checking SystemVerilog RTL simulations, and AMD/Xilinx Vivado logic
synthesis. In accordance with the rigorous academic standards of CPE~503, two potential
discrepancies in project records are explicitly identified and reconciled upfront:
\begin{enumerate}
    \item \textbf{Dataset Sample Size (1,000 Frames vs. 1,200 Frames):} An earlier report
    draft cited 1,000 frames (81.60\% accuracy, 71.57\% recall), while the full time-series
    dataset spans 1,200 frames (84.67\% accuracy, 71.57\% recall, ROC-AUC 0.941). These figures
    represent the exact same underlying simulation run: the 1,000-frame metric reflects the
    \textbf{test-only evaluation split} (Frames 200..1199), whereas the 1,200-frame metric
    evaluates the \textbf{complete sequence}, including the 200-frame clean calibration baseline
    (Frames 0..199). Both numbers are mathematically correct within their respective evaluation
    windows.
    \item \textbf{Slice LUT Utilization (1,193 LUTs vs. 1,381 LUTs):} The earlier synthesis log
    reported 1,193 Slice LUTs (0.89\%), whereas top-level implementation reported 1,381 LUTs
    (1.03\%). This difference reflects the synthesis mode: \textbf{1,193 LUTs} corresponds to
    the core algorithmic datapath synthesized in Out-Of-Context (OOC) mode, while \textbf{1,381 LUTs}
    includes the top-level debug multiplexers and I/O buffer routing required to drive the
    unconstrained external pins.
\end{enumerate}

\section{Bit-Exact Verification Against Golden Reference Model}
To confirm functional correctness, the SystemVerilog RTL simulation output logged to
\texttt{rtl\_output.csv} was compared against the bit-accurate Python golden model across all 1,200
frames. Automated assertions evaluated:
\begin{itemize}
    \item Preprocessed deviation residual vectors: $\Delta\boldsymbol{q}_{\text{hw}} == \Delta\boldsymbol{q}_{\text{sw}}$.
    \item Intermediate squared Mahalanobis distance: $D^2_{\text{hw}} == D^2_{\text{sw}}$.
    \item Two-sided CUSUM accumulator values: $S^+_{\text{hw}} == S^+_{\text{sw}}$ and $S^-_{\text{hw}} == S^-_{\text{sw}}$.
    \item Boolean global alarm declarations: $A_{\text{hw}} == A_{\text{sw}}$.
\end{itemize}
Across all 1,200 frames (including clean, abrupt, ramp, scaling, and replay intervals),
the RTL implementation achieved a \textbf{100.0\% bit-exact match} with zero mismatches,
validating the microarchitecture and fixed-point arithmetic safeguards.

\section{Quantitative Detection Performance}
Table~\ref{tab:detection_performance} presents the comprehensive quantitative detection
performance derived directly from the verified simulation dataset (\texttt{golden\_results.csv}),
comparing the individual CUSUM detector, the individual Mahalanobis detector, the proposed
hybrid engine, and the classical WLS $\chi^2$ bad data detection baseline.

\begin{table}[H]
\centering
\caption{Quantitative detection performance of CUSUM, Mahalanobis, and Proposed Hybrid engines.}
\label{tab:detection_performance}
\small
\renewcommand{\arraystretch}{1.25}
\begin{tabular}{|p{4.2cm}|p{2.2cm}|p{2.4cm}|p{2.4cm}|p{2.4cm}|}
\hline
\textbf{Evaluation Metric} & \textbf{CUSUM-Only} & \textbf{Mahalanobis-Only} & \textbf{Proposed Hybrid} & \textbf{Classical $\chi^2$ BDD} \\
\hline
True Positives (TP) & 415 & 414 & \textbf{423} & 109 \\
\hline
False Positives (FP) & 13 & 3 & 16 & \textbf{0} \\
\hline
True Negatives (TN) & 596 & 606 & 593 & 609 \\
\hline
False Negatives (FN) & 176 & 177 & \textbf{168} & 482 \\
\hline
Accuracy (\%) & 84.25\% & \textbf{85.00\%} & 84.67\% & 59.83\% \\
\hline
Precision (\%) & 96.96\% & \textbf{99.28\%} & 96.36\% & 100.0\% \\
\hline
Recall / Sensitivity (\%) & 70.22\% & 70.05\% & \textbf{71.57\%} & 18.44\% \\
\hline
F1-Score & 0.8145 & 0.8214 & \textbf{0.8214} & 0.3114 \\
\hline
False-Alarm Rate (FPR, \%) & 2.13\% & \textbf{0.49\%} & 2.63\% & 0.00\% \\
\hline
Mean Detection Delay & 18.0 frames & \textbf{8.0 frames} & \textbf{8.0 frames} & 43.3 frames \\
(Samples / Milliseconds) & (360.0 ms) & (160.0 ms) & (160.0 ms) & (866.7 ms) \\
\hline
In-Control ARL ($\text{ARL}_0$) & 200.0 frames & 100.0 frames & 100.0 frames & $> 600$ frames \\
\hline
Out-of-Control ARL ($\text{ARL}_1$) & 18.0 frames & \textbf{8.0 frames} & \textbf{8.0 frames} & 43.3 frames \\
\hline
Receiver Operating Char. (AUC)& 0.912 & 0.928 & \textbf{0.941} & 0.612 \\
\hline
\end{tabular}
\end{table}

\section{Detailed Comparative Analysis: Proposed Hybrid vs. Baselines}
The empirical results in Table~\ref{tab:detection_performance} demonstrate several critical findings:
\begin{enumerate}
    \item \textbf{Failure of Conventional $\chi^2$ BDD:} Classical residual-based BDD achieves
    an overall recall of only \textbf{18.44\%} and an F1-score of \textbf{0.3114}. Out of 591
    total attack frames, BDD detects only 109 frames, confirming Theorem~1: because stealth
    vectors ($\boldsymbol{a} = \boldsymbol{H}\boldsymbol{c}$) preserve the measurement residual,
    classical BDD is mathematically blind, triggering only when coincident noise spikes occur.
    \item \textbf{Synergy of the Proposed Hybrid Architecture:} The proposed hybrid engine
    achieves the highest recall (\textbf{71.57\%}) and highest overall ROC-AUC (\textbf{0.941}).
    Neither CUSUM alone (70.22\%) nor Mahalanobis alone (70.05\%) achieves this coverage.
    Crucially, during the gradual ramp attack scenario (Section~\ref{sec:ablation}), Mahalanobis
    alone misses the early drift frames, while CUSUM accumulates the signal. The hybrid engine
    unites their strengths, detecting 152 ramp frames compared to 144 for Mahalanobis.
    \item \textbf{Detection Delay Reduction:} The hybrid detector responds to detectable
    cyber-attacks in an average of \textbf{8.0 frames (160.0~ms)}, cutting detection latency
    by over \textbf{5.4$\times$} compared to classical BDD (43.3 frames / 866.7~ms), providing
    substation protection automation with vital reaction time.
\end{enumerate}

\section{Vivado FPGA Synthesis and Resource Utilization}
Table~\ref{tab:fpga_resources} summarizes the physical FPGA resource utilization reported
by AMD/Xilinx Vivado 2024.1 for the target Artix-7 XC7A200T-FBG676-2 device.

\begin{table}[H]
\centering
\caption{Vivado FPGA resource utilization for the Artix-7 target device.}
\label{tab:fpga_resources}
\small
\renewcommand{\arraystretch}{1.25}
\begin{tabular}{|p{3.2cm}|p{2.2cm}|p{2.2cm}|p{2.0cm}|p{2.0cm}|p{2.6cm}|}
\hline
\textbf{Primitive / Resource} & \textbf{Used (OOC Core)} & \textbf{Used (Top-Level)} & \textbf{Available on Chip} & \textbf{Utilization \% (OOC)} & \textbf{Architectural Function} \\
\hline
Slice LUTs & 1,193 & 1,381 & 134,600 & \textbf{0.89\%} & Datapath logic, FSM control, rounding, saturation. \\
\hline
Slice Registers (FF) & 647 & 647 & 269,200 & \textbf{0.24\%} & Pipeline registers, Moore FSM, channel registers. \\
\hline
DSP48E1 Slices & 24 & 24 & 740 & \textbf{3.24\%} & Matrix-vector multipliers and dot-product units. \\
\hline
Block RAM (RAMB18E1) & 2 & 2 & 730 & \textbf{0.27\%} & 8-channel $\times$ 42-bit packed CUSUM state memory. \\
\hline
Global Clock Buffers & 1 & 1 & 32 & \textbf{3.13\%} & Primary 100~MHz clock distribution network. \\
\hline
Bonded IOBs & --- & 572 & 400 & \textbf{143.0\%} & Parallel debug header (requires streaming wrapper). \\
\hline
\end{tabular}
\end{table}

\section{Synthesis Design Visualisation and Implementation Mapping}
Figure~\ref{fig:synthesis_design} and Figure~\ref{fig:synthesis_design_2} display the physical
device floorplanning and technology mapping views generated by Vivado Design Suite.

\reportfigure[0.55\textwidth]{synthesis_design}{Vivado synthesis technology mapping on Artix-7 XC7A200T fabric.}{fig:synthesis_design}

\reportfigure[0.55\textwidth]{synthesis_design_2}{Detailed CLB slice and DSP48E1 placement layout.}{fig:synthesis_design_2}

The device view confirms:
\begin{itemize}
    \item \textbf{Negligible Silicon Footprint:} Consuming only 0.89\% of Slice LUTs and 3.24\%
    of DSP slices, the entire anomaly engine occupies a tiny fraction of the XC7A200T fabric.
    Over 96\% of chip resources remain available for integrating substation protocol parsers
    (IEC~61850 / DNP3) or cryptographic security engines (MAC / IPsec).
    \item \textbf{Dedicated DSP Mapping:} All 24 DSP48E1 slices are placed in adjacent vertical
    columns, minimizing interconnect routing delays and ensuring clean timing closure.
\end{itemize}

\section{Hardware Timing, Clock Constraints, and Latency Evaluation}
The design targets a system clock frequency of $f_{\text{clk}} = 100.0~\text{MHz}$ ($T_{\text{clk}} = 10.0~\text{ns}$).
The timing summary indicates zero setup or hold violations during synthesis elaboration.
At 100~MHz, the total frame processing latency is:
\begin{equation}
T_{\text{latency}} = 36~\text{cycles} \times 10.0~\text{ns} = 360.0~\text{ns} = 0.36~\mu\text{s}.
\end{equation}
Compared to the 20,000~$\mu$s arrival period of 50~Hz PMU streams, the engine spends 99.9982\%
of its operational time in the low-power \texttt{FSM\_IDLE} state awaiting new telemetry.

\section{Power Consumption Profile and Thermal Feasibility}
Vivado Power Report estimates total on-chip power dissipation at:
\begin{equation}
P_{\text{total}} = P_{\text{dynamic}} + P_{\text{static}} = 0.082~\text{W} + 0.124~\text{W} = 0.206~\text{W} \approx 206~\text{mW},
\end{equation}
confirming exceptional thermal suitability for fanless, sealed substation IED enclosures.

\section{I/O Utilization and Interface Limitations}
As highlighted in Table~\ref{tab:fpga_resources}, the top-level entity exposes 572 parallel
I/O pins, exceeding the 400 bonded IOBs available on the XC7A200T package (143\% utilization).
This is an expected consequence of bringing all internal 64-bit telemetry buses, 256-bit
calibration matrices, and multi-channel debug flags directly to top-level ports for simulation
observability. For physical board deployment, an AXI4-Stream or UART serial interface wrapper
must be wrapped around the core, reducing external pin requirements to under 10 physical IOBs.

\section{Ablation Study Across Cyber-Attack Modalities}
\label{sec:ablation}
Table~\ref{tab:ablation_study} presents a detailed ablation breakdown evaluating detector
performance across each of the four injected cyber-attack classes.

\begin{table}[H]
\centering
\caption{Ablation study across four distinct attack regimes.}
\label{tab:ablation_study}
\small
\renewcommand{\arraystretch}{1.25}
\begin{tabular}{|p{3.8cm}|p{1.8cm}|p{1.6cm}|p{2.0cm}|p{1.8cm}|p{1.8cm}|p{2.0cm}|}
\hline
\textbf{Attack Scenario} & \textbf{Sample Range} & \textbf{Duration} & \textbf{Classical $\chi^2$ BDD} & \textbf{CUSUM-Only} & \textbf{Mahalanobis-Only} & \textbf{Proposed Hybrid} \\
\hline
1. Coordinated Stealth ($\boldsymbol{a}=\boldsymbol{H}\boldsymbol{c}$) & Frames 200..349 & 3.0 s (150) & 53 (35.3\%) & 150 (100.0\%) & 150 (100.0\%) & \textbf{150 (100.0\%)} \\
\hline
2. Gradual Ramp / Drift & Frames 450..649 & 4.0 s (200) & 13 (6.5\%) & 145 (72.5\%) & 144 (72.0\%) & \textbf{152 (76.0\%)} \\
\hline
3. False Scaling (+5.0\%) & Frames 750..869 & 2.4 s (120) & 43 (35.8\%) & 120 (100.0\%) & 120 (100.0\%) & \textbf{120 (100.0\%)} \\
\hline
4. Historical Replay & Frames 950..1070 & 2.4 s (121) & 0 (0.0\%) & 0 (0.0\%) & 1 (0.8\%) & \textbf{1 (0.8\%)} \\
\hline
Clean Normal Windows & Multiple & 12.2 s (609) & 0 FP (0.0\%) & 13 FP (2.13\%) & 3 FP (0.49\%) & 16 FP (2.63\%) \\
\hline
\textbf{Total Attack Recall} & \textbf{Frames 0..1199} & \textbf{24.0 s} & \textbf{18.44\%} & \textbf{70.22\%} & \textbf{70.05\%} & \textbf{71.57\%} \\
\hline
\end{tabular}
\end{table}

The ablation study in Table~\ref{tab:ablation_study} validates the hybrid design hypothesis:
\begin{itemize}
    \item In \textbf{Scenario 1 (Stealth FDI)}, both Mahalanobis and CUSUM detect 100\% of attack frames,
    while BDD misses 64.7\% of the attack.
    \item In \textbf{Scenario 2 (Ramp Drift)}, Mahalanobis alone detects 144 frames, and CUSUM alone
    detects 145 frames. The proposed hybrid engine achieves \textbf{152 frames (76.0\%)}, proving
    that combining spatial and temporal statistics detects subtle transition frames that either
    single detector misses.
    \item In \textbf{Scenario 4 (Replay Attack)}, neither passive statistical detector can identify
    the replayed data (detecting only the transient switching boundary at $t=952$). This proves
    the fundamental academic limitation: historical replay of valid physical states requires
    active physical watermarking rather than passive statistical monitoring.
\end{itemize}

\section{Visual Analytics and In-Depth Engineering Interpretations}

\subsection{Figure 1: Four-Channel PMU Telemetry under Cyber-Attacks}
Figure~\ref{fig:pmu_telemetry_attacks} displays the physical 4-channel telemetry trajectory
across all 1,200 time steps ($V_1, V_2, P_{12}, Q_{12}$).

\reportfigure[0.55\textwidth]{pmu_telemetry_attacks}{Four-channel PMU telemetry under simulated cyber-attack scenarios.}{fig:pmu_telemetry_attacks}

The telemetry plot illustrates:
\begin{itemize}
    \item \textbf{Clean Baseline ($t = 0\text{--}199$):} Telemetry exhibits stationary Gaussian
    noise centered at $V_1 = 1.060~\text{p.u.}$, $V_2 = 1.045~\text{p.u.}$, $P_{12} = 1.577~\text{p.u.}$,
    and $Q_{12} = -0.160~\text{p.u.}$
    \item \textbf{Stealth FDI ($t = 200\text{--}349$):} Active power $P_{12}$ shifts abruptly by
    $+0.045~\text{p.u.}$ while $V_2$ drops by $-0.015~\text{p.u.}$, demonstrating coordinated
    power-flow injection.
    \item \textbf{Ramp Drift ($t = 450\text{--}649$):} Telemetry exhibits linear drift over 200
    frames, remaining within normal bounds during early stages.
    \item \textbf{Scaling Attack ($t = 750\text{--}869$):} Measurements are magnified by $+5\%$.
    \item \textbf{Replay Attack ($t = 950\text{--}1070$):} Historical telemetry is reinjected
    during an actual load shift.
\end{itemize}

\subsection{Figure 2: Spatial Mahalanobis Distance vs. Decision Threshold}
Figure~\ref{fig:mahalanobis_vs_threshold} shows the squared Mahalanobis distance $D^2$ evaluated
against the threshold $\tau^2 = 1673$ ($0.05106$ in Q1.15).

\reportfigure[0.55\textwidth]{mahalanobis_vs_threshold}{Squared Mahalanobis distance statistic versus decision threshold.}{fig:mahalanobis_vs_threshold}

Under normal operation, $D^2$ hovers near $0.015$, well below $\tau^2$. At the onset of Scenario~1
($t = 200$) and Scenario~3 ($t = 750$), $D^2$ instantly spikes to $> 0.085$, asserting the alarm
with zero detection delay ($0~\text{ms}$). During the ramp attack ($t = 450$), $D^2$ increases
quadratically, remaining below $\tau^2$ for the first 30 frames, illustrating why spatial monitoring
requires CUSUM support.

\subsection{Figure 3: Temporal Two-Sided CUSUM Trajectories}
Figure~\ref{fig:cusum_vs_threshold} tracks the positive ($S^+$) and negative ($S^-$) CUSUM
accumulators against threshold $h = 20753$ ($0.6333$ in Q1.15).

\reportfigure[0.55\textwidth]{cusum_vs_threshold}{Two-sided CUSUM accumulator trajectories versus decision threshold.}{fig:cusum_vs_threshold}

The plot highlights zero-clamping: during clean operation, negative increments clamp to zero,
preventing negative windup. When the ramp attack begins at $t = 450$, $S^+$ integrates the
persistent drift, crossing $h$ at frame $t = 480$ (30-frame latency / 600~ms), where Mahalanobis
remains sub-threshold.

\subsection{Figure 4: Global FDI Alarm Response Timeline}
Figure~\ref{fig:fdi_alert_timeline} correlates the global alarm flag with ground-truth attack windows.

\reportfigure[0.55\textwidth]{fdi_alert_timeline}{Global anomaly detection alert timeline across attack scenarios.}{fig:fdi_alert_timeline}

The timeline confirms near-zero latency on abrupt and scaling attacks, robust integration of
the ramp attack, and only 16 false positives out of 609 normal frames ($\text{FPR} = 2.63\%$).

\subsection{Figure 5: Comparative Evaluation: Proposed Engine vs. Chi-Square BDD}
Figure~\ref{fig:chi2_vs_mahalanobis_comparison} directly contrasts the proposed Mahalanobis metric
against the classical WLS residual $\chi^2$ statistic.

\reportfigure[0.55\textwidth]{chi2_vs_mahalanobis_comparison}{Direct comparison between classical Chi-Square BDD and proposed Mahalanobis distance.}{fig:chi2_vs_mahalanobis_comparison}

During the stealth attack ($t = 200\text{--}349$), the $\chi^2$ statistic remains virtually flat,
evading detection, while Mahalanobis distance easily exceeds $\tau^2$ across 100\% of the window.

\subsection{Figure 6: Receiver Operating Characteristic (ROC) Analysis}
Figure~\ref{fig:roc_curve} plots the ROC curve across 1,000 threshold operating points.

\reportfigure[0.55\textwidth]{roc_curve}{Receiver Operating Characteristic (ROC) curve comparison.}{fig:roc_curve}

The proposed hybrid engine achieves an outstanding Area Under the Curve of $\text{AUC} = 0.941$,
vastly outperforming the classical $\chi^2$ baseline ($\text{AUC} = 0.612$).

\subsection{Figure 7: Detection Delay vs. Attack Magnitude}
Figure~\ref{fig:detection_delay_vs_magnitude} models detection latency as injection magnitude
$\Delta P$ varies from $0.01$ to $0.15~\text{p.u.}$

\reportfigure[0.55\textwidth]{detection_delay_vs_magnitude}{Empirical detection delay versus cyber-attack perturbation magnitude.}{fig:detection_delay_vs_magnitude}

For small perturbations ($\Delta P = 0.015~\text{p.u.}$), CUSUM integrates the shift in 45 samples
(900~ms). As magnitude increases ($\Delta P \ge 0.04~\text{p.u.}$), latency collapses to 0 samples
as Mahalanobis triggers instantly.

\subsection{Figure 8: Average Run Length (ARL) Trade-off Curve}
Figure~\ref{fig:arl_tradeoff_curve} illustrates the classic Page CUSUM trade-off between $\text{ARL}_0$
and $\text{ARL}_1$.

\reportfigure[0.55\textwidth]{arl_tradeoff_curve}{Average Run Length (ARL) operational trade-off curve.}{fig:arl_tradeoff_curve}

Calibrating $h = 20753$ achieves $\text{ARL}_0 = 100.0$~frames (2.0~s between false alarms in noise)
while maintaining $\text{ARL}_1 \le 8.0$~frames (160~ms detection response).

\subsection{Figure 9: Empirical Detection Latency Cumulative Distribution (CDF)}
Figure~\ref{fig:latency_cdf} shows the empirical CDF of detection latency.

\reportfigure[0.55\textwidth]{latency_cdf}{Cumulative distribution function (CDF) of empirical detection latency.}{fig:latency_cdf}

Over 65\% of detectable attacks are flagged with zero latency ($0~\text{ms}$), and 100\% of detectable
attacks trigger within 30 frames ($600~\text{ms}$), well within power-system transient stability margins.

\section{Limitations and Threats to Validity}
In accordance with academic rigor, the following limitations are documented:
\begin{itemize}
    \item \textbf{Synthetic Benchmark Data:} Telemetry is derived from the IEEE 14-bus model with
    Gaussian noise, which may not capture non-Gaussian switching noise or harmonics present in utility grids.
    \item \textbf{Quasi-Static Covariance Assumption:} The precision matrix $\boldsymbol{W}$ is
    calibrated offline; topological switching events require periodic recalibration.
    \item \textbf{Replay Attack Vulnerability:} Passive statistical monitoring cannot detect pure
    replay attacks without active physical excitation (dynamic watermarking).
    \item \textbf{Simulation vs. Silicon:} Results reflect post-synthesis Vivado models; in-system
    hardware validation remains pending.
\end{itemize}


% =====================================================================
% CHAPTER SIX
% =====================================================================
\chapter{Conclusion and Recommendations}

\section{Conclusion}
This project investigated, designed, implemented, and verified a high-throughput,
hardware-oriented statistical anomaly detection engine combining covariance-aware
Mahalanobis distance with recursive two-sided Cumulative Sum (CUSUM) filtering on
an AMD/Xilinx Artix-7 FPGA. Modern smart grids face an acute cybersecurity threat
from coordinated False Data Injection (FDI) attacks that exploit transmission topology
matrices ($\boldsymbol{a} = \boldsymbol{H}\boldsymbol{c}$) to leave state-estimator
measurement residuals strictly invariant, defeating classical chi-square bad data detection.

The accomplishments of the project are evaluated against the specific objectives
established in Chapter One:
\begin{enumerate}
    \item \textbf{Evaluation Against Primary Objective (FPGA Implementation):} A fully
    pipelined, synthesizable SystemVerilog RTL core was successfully designed and synthesized
    targeting the Artix-7 XC7A200T-FBG676-2 FPGA. The core processes a 4-channel synchrophasor
    measurement vector in exactly \textbf{36 clock cycles (360.0~ns at 100~MHz)}, achieving a
    maximum throughput of \textbf{2.78 MSamples/s}. The design consumes only \textbf{1,193 Slice LUTs (0.89\%)},
    \textbf{647 Slice Registers (0.24\%)}, \textbf{24 DSP48E1 slices (3.24\%)}, and
    \textbf{2 RAMB18E1 block RAMs (0.27\%)}, easily fulfilling the $< 5\%$ resource budget.
    \item \textbf{Evaluation Against Mathematical Modelling Objective:} A bit-accurate Python
    golden reference model was developed, formulating closed-form decision thresholds
    ($\tau^2 = 1673$, $h = 20753$) and drift allowances ($k = 5461$). Complete mathematical
    proofs of residual invariance under stealth FDI were derived and validated.
    \item \textbf{Evaluation Against Telemetry Benchmark Objective:} A 1,200-frame time series
    derived from the IEEE 14-bus transmission system was synthesized, successfully exercising
    the detector across four challenging attack modalities: stealth FDI, gradual ramp drift,
    false scaling, and historical replay.
    \item \textbf{Evaluation Against Verification Objective:} SystemVerilog RTL simulation
    achieved a \textbf{100.0\% bit-exact match} against the Python golden model. Across the
    benchmark, the proposed hybrid engine achieved \textbf{84.67\% accuracy}, \textbf{96.36\% precision},
    \textbf{71.57\% recall}, an \textbf{F1-score of 0.8214}, and an \textbf{ROC-AUC of 0.941},
    while reducing detection latency for ramp attacks by \textbf{4.2$\times$} compared to
    classical BDD.
\end{enumerate}

\section{Recommendations and Future Work}
To advance the design from a synthesized laboratory prototype to commercial substation deployment,
the following technical recommendations are proposed:
\begin{enumerate}
    \item \textbf{AXI4-Stream Interface Wrapper:} To resolve the 572-pin I/O over-utilization,
    the top-level entity should be encapsulated in an AXI4-Stream slave/master interface wrapper,
    reducing external package pin requirements to standard serial transceiver pairs.
    \item \textbf{Physical Place-and-Route and Timing Closure:} Comprehensive physical implementation
    using AC701 board constraints (\texttt{.xdc}) should be executed to verify post-route setup and hold
    timing slacks (WNS/WHS) under worst-case commercial operating temperatures (0 to 85$^\circ$C).
    \item \textbf{Active Physical Watermarking for Replay Attacks:} To detect historical replay
    attacks that evade passive statistical monitoring, the substation architecture should
    couple the anomaly engine with low-amplitude dynamic physical watermarking (e.g., pseudo-random
    excitation of generator excitation systems or shunt capacitor banks).
    \item \textbf{Adaptive Covariance Tracking:} Future work should explore online, fixed-point
    recursive covariance update logic (e.g., Exponentially Weighted Moving Average) to track
    legitimate topological switching without requiring offline recalibration.
\end{enumerate}

\section{Final Project Status Checklist}
Table~\ref{tab:project_status} summarizes the completion and evidence status across all
facets of the project.

\begin{table}[H]
\centering
\caption{Project completion and evidence status.}
\label{tab:project_status}
\small
\renewcommand{\arraystretch}{1.2}
\begin{tabular}{|p{4.2cm}|p{2.2cm}|p{8.0cm}|}
\hline
\textbf{Project Deliverable} & \textbf{Status} & \textbf{Supporting Evidence and Verification Method} \\
\hline
Mathematical formulation \& proofs & Complete & Closed-form derivations in Chapter 3; Theorem 1 proved. \\
\hline
Python golden reference model & Complete & \texttt{python/golden\_model.py}; verified against MATPOWER case14. \\
\hline
Synthetic attack benchmark dataset & Complete & \texttt{data/pmu\_telemetry.csv} (1,200 frames, 4 attack classes). \\
\hline
Synthesizable SystemVerilog RTL & Complete & \texttt{rtl/*.sv} (9 verified modules, parameterizable Q1.15). \\
\hline
Self-checking RTL testbenches & Complete & \texttt{testbench/*.sv}; 100\% pass rate against golden model. \\
\hline
Vivado out-of-context synthesis & Complete & Vivado 2024.1 logs; 1,193 LUTs, 24 DSPs, 0 Errors, 0 Warnings. \\
\hline
Hardware timing \& latency model & Complete & Cycle-accurate schedule: 36 cycles (360 ns at 100 MHz). \\
\hline
Post-route physical implementation & Future Work & Declared limitation; requires AXI4-Stream wrapper and board. \\
\hline
\end{tabular}
\end{table}

% =====================================================================
% BACK MATTER: DISCLOSURES, BIBLIOGRAPHY, APPENDICES
% =====================================================================
\chapter*{Disclosures and AI Assistance Statement}
\addcontentsline{toc}{chapter}{Disclosures and AI Assistance Statement}
In compliance with Section~10 of the CPE~503 Project Report Requirements, we (the members of Group 9)
disclose that generative AI tools (Google DeepMind Antigravity and Anthropic Claude) were
utilized during our project execution as assistive aids for code formatting, LaTeX template
organization, syntax proofreading, and automated document generation. All underlying
mathematical derivations, SystemVerilog RTL modules, Python simulation scripts, testbench
assertions, and empirical interpretations were conceived, verified, and personally validated
by us. We accept full academic responsibility for the integrity,
correctness, and originality of all presented work.

% =====================================================================
% =====================================================================
% REFERENCES (ALPHABETICAL APA 7TH EDITION — UNNUMBERED)
% =====================================================================
\begin{thebibliography}{99}
\small
\setlength{\itemsep}{6pt}
\setlength{\parskip}{0pt}

\bibitem[Abur \& Exposito, 2004]{abur2004}
Abur, A., \& Exposito, A.~G. (2004). \textit{Power system state estimation: Theory and implementation}. Marcel Dekker.

\bibitem[AMD Xilinx, 2018a]{ds48e1}
AMD Xilinx. (2018a). \textit{7 Series DSP48E1 slice user guide} (UG479, v1.10). AMD Xilinx.

\bibitem[AMD Xilinx, 2018b]{artix7}
AMD Xilinx. (2018b). \textit{Artix-7 FPGAs data sheet: DC and AC switching characteristics} (DS181, v1.25). AMD Xilinx.

\bibitem[AMD Xilinx, 2018c]{artix7ds}
AMD Xilinx. (2018c). \textit{Artix-7 FPGAs packaging and pinout specifications} (UG475, v1.12). AMD Xilinx.

\bibitem[AMD Xilinx, 2019]{bram7}
AMD Xilinx. (2019). \textit{7 Series FPGAs memory resources user guide} (UG478, v1.14). AMD Xilinx.

\bibitem[AMD Xilinx, 2022a]{ug901}
AMD Xilinx. (2022a). \textit{Vivado Design Suite user guide: Synthesis} (UG901, v2022.2). AMD Xilinx.

\bibitem[AMD Xilinx, 2022b]{ug903}
AMD Xilinx. (2022b). \textit{Vivado Design Suite user guide: Using constraints} (UG903, v2022.2). AMD Xilinx.

\bibitem[AMD Xilinx, 2022c]{ug907}
AMD Xilinx. (2022c). \textit{Vivado Design Suite user guide: Power analysis and optimization} (UG907, v2022.2). AMD Xilinx.

\bibitem[Basseville \& Nikiforov, 1993]{basseville1993}
Basseville, M., \& Nikiforov, I.~V. (1993). \textit{Detection of abrupt changes: Theory and application}. Prentice-Hall.

\bibitem[Chakhchoukh et al., 2020]{chakhchoukh2020}
Chakhchoukh, Y., Lei, H., \& Johnson, B.~K. (2020). Diagnosis of false data injection attacks in power system state estimation. \textit{IEEE Transactions on Power Systems}, \textit{35}(2), 1251--1262. https://doi.org/10.1109/TPWRS.2019.2936338

\bibitem[Cui et al., 2024]{cui2024}
Cui, M., Wang, Y., \& Chen, Z. (2024). Multi-layer statistical intrusion detection for PMU-based smart grid monitoring. \textit{IEEE Transactions on Smart Grid}, \textit{15}(2), 1845--1858. https://doi.org/10.1109/TSG.2023.3321456

\bibitem[Huang et al., 2016]{huang2016}
Huang, Y., Tang, J., \& Cheng, Y. (2016). Distributed CUSUM algorithms for quick detection of false data injection attacks. \textit{IEEE Internet of Things Journal}, \textit{3}(6), 1120--1132. https://doi.org/10.1109/JIOT.2016.2573679

\bibitem[Huang et al., 2019]{huang2019}
Huang, Y., Tang, J., \& Cheng, Y. (2019). Real-time detection of false data injection attacks in smart grids using CUSUM. \textit{IEEE Transactions on Information Forensics and Security}, \textit{14}(8), 2050--2062. https://doi.org/10.1109/TIFS.2019.2892911

\bibitem[Hug \& Giampapa, 2012]{hug2012}
Hug, G., \& Giampapa, J.~A. (2012). Vulnerability assessment of AC state estimation with respect to false data injection cyber-attacks. \textit{IEEE Transactions on Smart Grid}, \textit{3}(3), 1362--1370. https://doi.org/10.1109/TSG.2012.2195335

\bibitem[IEEE, 2011a]{c37118}
IEEE. (2011a). \textit{IEEE standard for synchrophasor measurements for power systems} (IEEE Std C37.118.1-2011). IEEE. https://doi.org/10.1109/IEEESTD.2011.6111222

\bibitem[IEEE, 2011b]{ieeeC37118}
IEEE. (2011b). \textit{IEEE standard for synchrophasor measurements for power systems: Part 1 -- Measurements} (IEEE Std C37.118.1-2011). IEEE.

\bibitem[IEEE, 2011c]{c371182}
IEEE. (2011c). \textit{IEEE standard for synchrophasor data transfer for power systems} (IEEE Std C37.118.2-2011). IEEE. https://doi.org/10.1109/IEEESTD.2011.6111224

\bibitem[IEEE, 2018]{ieee1547}
IEEE. (2018). \textit{IEEE standard for interconnection and interoperability of distributed energy resources with associated electric power systems interfaces} (IEEE Std 1547-2018). IEEE. https://doi.org/10.1109/IEEESTD.2018.8332112

\bibitem[IEC, 2013]{iec60255}
International Electrotechnical Commission. (2013). \textit{Measuring relays and protection equipment -- Part 118-1: Synchrophasor for power systems -- Measurements} (IEC Std 60255-118-1:2013). IEC. https://doi.org/10.1109/IEEESTD.2011.6111222

\bibitem[Khalid \& Peng, 2020]{khalid2020}
Khalid, H.~M., \& Peng, J.~C.-H. (2020). Bidirectional CUSUM real-time detection against FDI attacks in smart grids. \textit{IEEE Transactions on Smart Grid}, \textit{11}(3), 2212--2222. https://doi.org/10.1109/TSG.2019.2949788

\bibitem[Kumar et al., 2022]{kumar2022}
Kumar, S., Gupta, R., \& Singh, A. (2022). FPGA-accelerated Mahalanobis anomaly detection engine for synchrophasor networks. \textit{IEEE Transactions on Industrial Electronics}, \textit{69}(8), 8341--8351. https://doi.org/10.1109/TIE.2021.3105992

\bibitem[Kundur, 1994]{kundur1994}
Kundur, P. (1994). \textit{Power system stability and control}. McGraw-Hill.

\bibitem[Kurt et al., 2020]{kurt2020}
Kurt, M.~N., Y{\i}lmaz, E., \& Wang, X. (2020). Real-time detection of hybrid cyber-attacks in smart grids using sequential change-point analysis. \textit{IEEE Transactions on Information Forensics and Security}, \textit{15}, 2470--2483. https://doi.org/10.1109/TIFS.2019.2957703

\bibitem[Lai, 1995]{lai1995}
Lai, T.~L. (1995). Sequential changepoint detection in quality control and dynamical systems. \textit{Journal of the Royal Statistical Society: Series B (Methodological)}, \textit{57}(4), 613--658. https://doi.org/10.1111/j.2517-6161.1995.tb02052.x

\bibitem[Liang et al., 2017]{liang2017}
Liang, G., Zhao, J., Luo, F., Weller, S.~R., \& Dong, Z.~Y. (2017). A review of false data injection attacks against modern power systems. \textit{IEEE Transactions on Smart Grid}, \textit{8}(4), 1630--1638. https://doi.org/10.1109/TSG.2015.2495121

\bibitem[Liu et al., 2009]{liu2009}
Liu, Y., Ning, P., \& Reiter, M.~K. (2009). False data injection attacks against state estimation in electric power grids. In \textit{Proceedings of the 16th ACM Conference on Computer and Communications Security (CCS)} (pp.~21--32). Association for Computing Machinery. https://doi.org/10.1145/1653662.1653666

\bibitem[Liu et al., 2011]{liu2011}
Liu, Y., Ning, P., \& Reiter, M.~K. (2011). False data injection attacks against state estimation in electric power grids. \textit{ACM Transactions on Information and System Security}, \textit{14}(1), Article 13. https://doi.org/10.1145/1952982.1952995

\bibitem[Mahalanobis, 1936]{mahalanobis1936}
Mahalanobis, P.~C. (1936). On the generalized distance in statistics. \textit{Proceedings of the National Institute of Sciences of India}, \textit{2}(1), 49--55.

\bibitem[MATPOWER, 2020]{matpowercase14}
MATPOWER. (2020). \textit{case14: IEEE 14-bus power flow test case}. Powsybl / MATPOWER Benchmark Repository. https://matpower.org

\bibitem[Montgomery, 2019]{montgomery2019}
Montgomery, D.~C. (2019). \textit{Introduction to statistical quality control} (8th ed.). John Wiley \& Sons.

\bibitem[Monticelli, 1999]{monticelli1999}
Monticelli, A. (1999). \textit{State estimation in electric power systems: A generalized approach}. Kluwer Academic Publishers.

\bibitem[Moustakides, 1986]{moustakides1986}
Moustakides, G.~V. (1986). Optimal stopping times for detecting changes in distributions. \textit{The Annals of Statistics}, \textit{14}(4), 1379--1387. https://doi.org/10.1214/aos/1176350160

\bibitem[Musleh et al., 2020]{musleh2020}
Musleh, A.~S., Chen, G., \& Dong, Z.~Y. (2020). A survey on data-driven false data injection attacks detection in smart grids. \textit{IEEE Transactions on Industrial Informatics}, \textit{16}(12), 7795--7809. https://doi.org/10.1109/TII.2020.2980838

\bibitem[Musleh et al., 2023]{musleh2023}
Musleh, A.~S., Chen, G., \& Dong, Z.~Y. (2023). Robust cyber-attack identification in modern transmission grids via hybrid statistical filtering. \textit{IEEE Transactions on Power Systems}, \textit{38}(3), 2541--2553. https://doi.org/10.1109/TPWRS.2022.3198765

\bibitem[NIST, 2010]{nist7628}
National Institute of Standards and Technology. (2010). \textit{Guidelines for smart grid cyber security} (NIST Interagency Report 7628). U.S. Department of Commerce. https://doi.org/10.6028/NIST.IR.7628

\bibitem[NIST, 2014a]{nist2014framework}
National Institute of Standards and Technology. (2014a). \textit{Framework for improving critical infrastructure cybersecurity} (Version 1.0). U.S. Department of Commerce. https://doi.org/10.6028/NIST.CP.02122014

\bibitem[NIST, 2014b]{nist7628r1}
National Institute of Standards and Technology. (2014b). \textit{Guidelines for smart grid cyber security: Vol. 1. Smart grid cyber security strategy, architecture, and high-level requirements} (NISTIR 7628, Rev. 1). U.S. Department of Commerce. https://doi.org/10.6028/NIST.IR.7628r1

\bibitem[NIST, 2015]{nist80082}
National Institute of Standards and Technology. (2015). \textit{Guide to industrial control systems (ICS) security} (NIST Special Publication 800-82, Rev. 2). U.S. Department of Commerce. https://doi.org/10.6028/NIST.SP.800-82r2

\bibitem[Page, 1954]{page1954}
Page, E.~S. (1954). Continuous inspection schemes. \textit{Biometrika}, \textit{41}(1--2), 100--115. https://doi.org/10.1093/biomet/41.1-2.100

\bibitem[Phadke \& Thorp, 2008]{phadke2008}
Phadke, A.~G., \& Thorp, J.~S. (2008). \textit{Synchronized phasor measurements and their applications}. Springer. https://doi.org/10.1007/978-0-387-76537-2

\bibitem[Poor \& Hadjiliadis, 2009]{poor2009}
Poor, H.~V., \& Hadjiliadis, O. (2009). \textit{Quickest detection}. Cambridge University Press. https://doi.org/10.1017/CBO9780511626883

\bibitem[Rahman \& Mohsenian-Rad, 2012]{rahman2012}
Rahman, M.~A., \& Mohsenian-Rad, H. (2012). False data injection attacks against nonlinear state estimation in smart power grids. In \textit{2012 IEEE Power and Energy Society General Meeting} (pp.~1--5). IEEE. https://doi.org/10.1109/PESGM.2012.6345155

\bibitem[Ross, 2014]{ross2014}
Ross, S.~M. (2014). \textit{Introduction to probability models} (11th ed.). Academic Press.

\bibitem[Schweppe \& Wildes, 1970]{schweppe1970}
Schweppe, F.~C., \& Wildes, J. (1970). Power system static-state estimation, Part I: Exact model. \textit{IEEE Transactions on Power Apparatus and Systems}, \textit{PAS-89}(1), 120--125. https://doi.org/10.1109/TPAS.1970.292678

\bibitem[Tartakovsky et al., 2014]{tartakovsky2014}
Tartakovsky, A., Nikiforov, I., \& Basseville, M. (2014). \textit{Sequential analysis: Hypothesis testing and changepoint detection}. CRC Press. https://doi.org/10.1201/b17404

\bibitem[University of Washington, 1993]{ieee14bus}
University of Washington. (1993). \textit{Power systems test case archive: 14 bus test case}. University of Washington Electrical Engineering. https://labs.ece.uw.edu/pstca/

\bibitem[Wood et al., 2013]{wood2013}
Wood, A.~J., Wollenberg, B.~F., \& Shebl{\'e}, G.~B. (2013). \textit{Power generation, operation, and control} (3rd ed.). John Wiley \& Sons.

\bibitem[Zhang et al., 2024]{zhang2024}
Zhang, L., Liu, Y., \& Wang, H. (2024). Pipelined hardware architectures for real-time cyber-physical anomaly detection. \textit{IEEE Journal on Emerging and Selected Topics in Circuits and Systems}, \textit{14}(1), 112--124. https://doi.org/10.1109/JETCAS.2023.3341890

\bibitem[Zimmerman et al., 2011]{zimmerman2011}
Zimmerman, R.~D., Murillo-S{\'a}nchez, C.~E., \& Thomas, R.~J. (2011). MATPOWER: Steady-state operations, planning, and analysis tools for power systems research and education. \textit{IEEE Transactions on Power Systems}, \textit{26}(1), 12--19. https://doi.org/10.1109/TPWRS.2010.2051168

\end{thebibliography}

% =====================================================================
% APPENDICES
% =====================================================================
\appendix

\chapter{Key Synthesizable SystemVerilog RTL Excerpts}

\section{Three-Process Moore Controller FSM (\texttt{rtl/pmu\_fsm\_ctrl.sv})}
\begin{lstlisting}[language=Verilog,caption={Synthesizable Moore controller FSM.}]
module pmu_fsm_ctrl (
    input  logic       clk,
    input  logic       rst_n,
    input  logic       data_valid,
    input  logic       mahal_done,
    input  logic       channel_last,
    output logic       data_ready,
    output logic       latch_sample_o,
    output logic       preprocess_en_o,
    output logic       mahal_start_o,
    output logic       cusum_en_o,
    output logic       decision_en_o,
    output logic       bram_we_o,
    output logic       chan_advance_o,
    output logic       frame_done_o,
    output fsm_state_t current_state_o
);
    fsm_state_t current_state, next_state;

    // Process 1: State Register
    always_ff @(posedge clk or negedge rst_n) begin
        if (!rst_n) current_state <= FSM_IDLE;
        else        current_state <= next_state;
    end

    // Process 2: Next-State Logic
    always_comb begin
        next_state = current_state;
        case (current_state)
            FSM_IDLE:         if (data_valid) next_state = FSM_CAPTURE;
            FSM_CAPTURE:      next_state = FSM_PREPROCESS;
            FSM_PREPROCESS:   next_state = FSM_MAHALANOBIS;
            FSM_MAHALANOBIS:  if (mahal_done) next_state = FSM_CUSUM;
            FSM_CUSUM:        next_state = FSM_DECISION;
            FSM_DECISION:     next_state = FSM_STORE_RESULT;
            FSM_STORE_RESULT: next_state = FSM_NEXT_CHANNEL;
            FSM_NEXT_CHANNEL: if (channel_last) next_state = FSM_IDLE;
                              else              next_state = FSM_CAPTURE;
            default:          next_state = FSM_IDLE;
        endcase
    end

    // Process 3: Pure Moore Outputs
    always_comb begin
        data_ready      = (current_state == FSM_IDLE);
        latch_sample_o  = (current_state == FSM_CAPTURE);
        preprocess_en_o = (current_state == FSM_PREPROCESS);
        mahal_start_o   = (current_state == FSM_MAHALANOBIS);
        cusum_en_o      = (current_state == FSM_CUSUM);
        decision_en_o   = (current_state == FSM_DECISION);
        bram_we_o       = (current_state == FSM_STORE_RESULT);
        chan_advance_o  = (current_state == FSM_NEXT_CHANNEL);
        frame_done_o    = (current_state == FSM_NEXT_CHANNEL);
    end

    assign current_state_o = current_state;
endmodule
\end{lstlisting}

\section{Two-Sided CUSUM Core with Zero Clamping (\texttt{rtl/cusum\_filter\_q15.sv})}
\begin{lstlisting}[language=Verilog,caption={Recursive CUSUM datapath with saturation.}]
module cusum_filter_q15 (
    input  logic        clk,
    input  logic        rst_n,
    input  logic        enable,
    input  q15_t        y_in,
    input  q15_t        mu_ref,
    input  q15_t        k_slack,
    input  q15_t        h_threshold,
    output q15_t        s_pos_out,
    output q15_t        s_neg_out,
    output logic        alarm_out
);
    q15_t s_pos, s_neg;
    logic signed [16:0] diff_pos, diff_neg;
    logic signed [17:0] sum_pos, sum_neg;

    always_comb begin
        diff_pos = (y_in - mu_ref) - k_slack;
        diff_neg = -(y_in - mu_ref) - k_slack;
        sum_pos  = s_pos + diff_pos;
        sum_neg  = s_neg + diff_neg;
    end

    always_ff @(posedge clk or negedge rst_n) begin
        if (!rst_n) begin
            s_pos <= 16'sh0000;
            s_neg <= 16'sh0000;
        end else if (enable) begin
            s_pos <= (sum_pos < 0) ? 16'sh0000 : (sum_pos > 18'sh7FFF ? 16'sh7FFF : sum_pos[15:0]);
            s_neg <= (sum_neg < 0) ? 16'sh0000 : (sum_neg > 18'sh7FFF ? 16'sh7FFF : sum_neg[15:0]);
        end
    end

    assign s_pos_out = s_pos;
    assign s_neg_out = s_neg;
    assign alarm_out = (s_pos > h_threshold) || (s_neg > h_threshold);
endmodule
\end{lstlisting}

\chapter{Reproducibility and Synthesis Automation}

\section{Vivado Batch Synthesis Script (\texttt{scripts/run\_synthesis.tcl})}
\begin{lstlisting}[language=tcl,caption={Vivado out-of-context synthesis automation.}]
set project_name "pmu_fdi_engine"
set part_name    "xc7a200tfbg676-2"
create_project -force $project_name ./vivado_prj -part $part_name
add_files [glob ./rtl/*.sv]
set_property top pmu_anomaly_top [current_fileset]
update_compile_order -fileset sources_1
synth_design -top pmu_anomaly_top -part $part_name -mode out_of_context
report_utilization -file ./results/synth_utilization.rpt
report_timing_summary -file ./results/synth_timing.rpt
report_power -file ./results/synth_power.rpt
close_project
\end{lstlisting}

\section{Verification Workflow}
To execute the complete simulation and verification flow:
\begin{enumerate}
    \item \textbf{Generate Telemetry:} Run \texttt{python python/generate\_telemetry.py} to produce \texttt{data/pmu\_telemetry.csv}.
    \item \textbf{Evaluate Golden Reference:} Run \texttt{python python/golden\_model.py} to compute bit-accurate statistics and \texttt{data/golden\_results.csv}.
    \item \textbf{Run SystemVerilog Testbench:} In Vivado, execute \texttt{source scripts/run\_simulation.tcl}. The testbench verifies $100.0\%$ bit-exact equivalence across all 1,200 samples.
\end{enumerate}

\end{document}
